\documentclass[usenatbib,twoside,twocolumn]{mnras}

% Paper packages
\usepackage{graphicx}
\usepackage{amsmath}
\usepackage{amssymb}
\usepackage{booktabs}

% Denser float packing (reduce float-only pages; no content change)
\renewcommand{\topfraction}{0.9}
\renewcommand{\bottomfraction}{0.85}
\renewcommand{\textfraction}{0.07}
\renewcommand{\floatpagefraction}{0.75}
\renewcommand{\dbltopfraction}{0.9}
\renewcommand{\dblfloatpagefraction}{0.75}
\setcounter{topnumber}{3}
\setcounter{bottomnumber}{2}
\setcounter{totalnumber}{5}
\setcounter{dbltopnumber}{3}

% Define astronomical units
\newcommand{\msun}{\ensuremath{M_{\odot}}}
\newcommand{\pc}{\ensuremath{\rm pc}}
\newcommand{\kms}{\ensuremath{\rm km\,s^{-1}}}

\title[Filament fragmentation and core spacing -- II]{Filament fragmentation and core spacing -- II.\\ Contracting filaments need not select the equilibrium wavelength}

\author[G.~J.~White]{G.~J.~White$^{1,2}$\\
$^{1}$School of Physical Sciences, The Open University, Walton Hall, Milton Keynes, MK7 6AA, UK\\
$^{2}$Rutherford Appleton Laboratory, Harwell Science and Innovation Campus, Didcot, OX11 0QX, UK}

\date{\today}
\pubyear{2026}

\begin{document}
\label{firstpage}
\maketitle

\begin{abstract}

Paper~I shows that dense cores along \textit{Herschel} filaments are clustered rather than periodically spaced. We ask what physics could produce such a distribution, beginning with the assumption on which the classical comparison rests: must a filament fragment at the fastest-growing wavelength $\lambda_{\rm m}=22H$ of the equilibrium cylinder from which it formed? We test that with ideal-MHD Athena++ calculations.

The answer rests on a control: a drag-relaxed Ostriker cylinder and a near-critical contracting filament pass through one identical growth-rate measurement. For the equilibrium the pipeline recovers the classical maximum at $22H$ to within $10$ per cent. For the contracting filament it resolves none: the growth rate rises monotonically to the resolution limit, leaving only the upper limit $\lambda<1.1\,\lambda_{\rm J}$. For the idealised configuration tested here, classical mode selection is not obligatory.

A ladder of impulsively contracted cylinders supports that reading without establishing a law: within a fixed transverse domain its resolved modes follow a gravitational $\rho^{-1/2}$ scaling, but enlarging the domain and extending the seeded modes breaks the rank ordering and does not recover it.

Every homogeneous channel fails. The calculations give a quasi-periodic bead chain, not the observed clustering; driven turbulence raises the fragment count while \emph{lowering} the gap dispersion and Fano factor; and a toroidally dominated helical equilibrium shrinks the width faster than the wavelength, so $\lambda/W$ rises. Only an axially varying line mass produces core-free stretches, because fragmentation is confined to the locally supercritical parts. That predicts $f_{\rm occ}\simeq f_{\rm elig}$, satisfied by Paper~I's two independent measurements in three clouds of four.

\end{abstract}

\begin{keywords}
stars: formation -- ISM: structure -- ISM: filaments -- ISM: magnetic fields -- MHD -- methods: numerical
\end{keywords}

\section{Introduction}
\label{sec:intro}

An isothermal, self-gravitating cylinder in hydrostatic equilibrium has a critical mass per unit length $M_{\rm line,crit}=2c_s^2/G$ \citep{Ostriker1964} and fragments most rapidly at $\lambda_{\rm m}=22H$, where $H=c_s(4\pi G\rho_c)^{-1/2}$ is the scale height \citep{Inutsuka1992,Nagasawa1987,Fischera2012}. In units of the observable width this is $\approx5$--$6\times$FWHM, and it has become the reference against which observed core separations are read.

The derivation assumes a stationary background. Real filaments are not stationary. They form out of, and continue to accrete from, a turbulent medium; their central densities rise while perturbations grow; and the observed population contains objects at a range of evolutionary stages \citep{Andre2014,Clarke2016,Clarke2017,Clarke2020}. Whether the equilibrium wavelength survives that is a question that can be answered by direct calculation, and this paper answers it.

The question is narrow and we state its limits at the outset. We ask: \emph{must a dynamically evolving filament select the fragmentation wavelength associated with its initial equilibrium?} We do not ask, and cannot answer with these calculations, whether dynamical evolution explains the core separations observed in real clouds. The companion paper shows that those separations are not described by a single wavelength at all, and that what a successful theory must reproduce is a spatial point process. Our calculations produce a periodic chain and therefore do not reproduce it. They are offered as a counterexample to a theoretical assumption, and nothing more. Table~\ref{tab:notation} fixes the notation.

\begin{table}
\caption{Notation. A \emph{limit} and a \emph{detection} are never written with the same symbol: $\lambda_{\rm turn}^{\rm lim}$ is an upper bound from a spectrum with no resolved maximum, $\lambda_{\rm max}$ a resolved one.}
\label{tab:notation}
\footnotesize
\setlength{\tabcolsep}{4pt}
\begin{tabular}{@{}lp{5.9cm}@{}}
\toprule
Symbol & Meaning \\
\midrule
$\lambda_{\rm J}$, $t_{\rm J}$ & Jeans length and time on the fixed ambient density \\
$H$ & scale height $c_s(4\pi G\rho_c)^{-1/2}$ on the instantaneous central density \\
$\lambda_{\rm m}=22H$ & fastest-growing wavelength of the equilibrium cylinder \\
$\Gamma(\lambda)$ & amplitude growth rate, $\tfrac12\,{\rm d}\ln P/{\rm d}t$ \\
$\lambda_{\rm max}$ & wavelength of a \emph{resolved, stable} maximum of $\Gamma$ \\
$\lambda_{\rm turn}^{\rm lim}$ & upper limit on any turnover where none is resolved \\
$\rho_{\rm eff}$ & central density at the mid-point of the linear growth window \\
$f$, $\beta$, $A$ & line-mass fraction, plasma beta, imposed inward-velocity amplitude \\
$W_{\rm core},W_{\rm form},W_{\rm fil}$ & the three widths of Table~\ref{tab:widthdefs} \\
$\eta_W=T_1$ & $W_{\rm form}/W_{\rm fil}$, the conversion to observable units \\
\bottomrule
\end{tabular}
\end{table}

\smallskip\noindent\textbf{Structure.} Section~\ref{sec:method} gives the numerical set-up and the initial conditions in full. Section~\ref{sec:control} is the principal result: the same measurement applied to a relaxed cylinder and to an evolving one. It is the only conclusion here that we regard as established independently of the transverse boundary condition. Section~\ref{sec:ladder} adds supporting, fixed-box evidence from a seventeen-run ladder in dynamical state (which we call the \emph{impulsive-contraction ladder}, to distinguish it from the line-mass ladder of Table~\ref{tab:production}) and extracts the scaling. Section~\ref{sec:convergence} sets out the resolution and convergence tests, including a $1024\times256^2$ check. Section~\ref{sec:fields} states what the campaign does and does not constrain about magnetic geometry, and Section~\ref{sec:ambient} what it establishes about ambient confinement. Section~\ref{sec:comparison} makes the comparison with the observations of Paper~I explicit about what it can and cannot mean.

\section{Method}
\label{sec:method}


\smallskip\noindent\emph{The calculations below test whether sub-equilibrium fragmentation scales are dynamically possible. They are not intended to reproduce the observed core point process, and no number quoted in them is a prediction of an observed spacing.}

\smallskip We use Athena++ \citep{Stone2020}, release of August 2021 with the \emph{isothermal} equation of state, $P=\rho c_s^2$. The energy equation is not evolved and no cooling prescription is applied; the evolved equations are continuity, momentum including the Lorentz force, and ideal induction, coupled to self-gravity through $\nabla^2\phi=4\pi G\rho$. We use the HLLD isothermal Riemann solver, piecewise-linear reconstruction in the primitive variables, the second-order van Leer integrator at CFL $0.3$, and constrained transport. Self-gravity is solved with the FFT Poisson solver, which is triply periodic; transverse isolation is therefore \emph{not} imposed exactly. We approximate it by placing the transverse boundary at $d\geq1\,\lambda_{\rm J}$ from the spine and subtracting the mean density, so that the periodic Green's function acts on a zero-mean source and the spurious transverse image forces fall as the images recede. The adequacy of that approximation is a property of the box, not of the code, and it is weaker than we would like. Reflecting transverse walls at $L_\perp/2=2\,\lambda_{\rm J}$ generate a two-dimensional lattice of image filaments at that spacing, whose transverse gravity opposes accretion onto the spine and therefore \emph{supports} it. Doubling $L_\perp$ to $8\,\lambda_{\rm J}$ at fixed cell size removes most of that support: the central density at the mid-point of the growth window rises from $1.42$ to $9.11\,\rho_0$ for the static rung of Section~\ref{sec:ladder} and from $1.48$ to $4.78\,\rho_0$ for the contracting rung at $A=0.6$. Both then select the shortest seeded mode, $n=10$, so both become upper limits and the doubled-box pair cannot be compared with the production pair on equal terms without extending the seeded mode set. Three consequences follow. First, the absolute value of $\rho_{\rm eff}$ in the production runs carries a transverse-box systematic that we have not bounded, so the normalisation $\lambda/22H(\rho_{\rm eff})$ is uncertain at the tens-of-per-cent level and not at the few-per-cent level. Second, comparisons made at the same box size preserve the qualitative ordering when the box is enlarged, the static and contracting rungs remaining clearly separated (Section~\ref{sec:ladder}); but that is a statement about the ordering, not about the slope or normalisation of the ladder relation, neither of which has been shown to be transverse-box converged. Third, we have run that larger box with an extended seeded mode set, and the outcome is reported in Section~\ref{sec:ladder}: the equilibrium-versus-contracting separation survives it, the fitted density scaling does not.

\smallskip\noindent\textbf{Units and conventions.} Lengths are in Jeans lengths, $\lambda_{\rm J}\equiv c_s(\pi/G\rho_{\rm bg})^{1/2}$, evaluated on the \emph{ambient} density $\rho_{\rm bg}$ and fixed for a given run; times in $t_{\rm J}\equiv(4\pi G\rho_{\rm bg})^{-1/2}$. The equilibrium scale height is $H\equiv c_s(4\pi G\rho_c)^{-1/2}$, evaluated on the instantaneous \emph{central} density, so $H$ evolves within a run and $\lambda_{\rm J}$ does not. Both appear because the classical result is stated in $H$ and the grid is specified in $\lambda_{\rm J}$; wherever a ratio is quoted we state which density it uses.

\smallskip\noindent\textbf{The initial condition.} The production runs use a Gaussian filament embedded in a uniform ambient medium,
\begin{equation}
\rho(r)=\rho_{\rm bg}+\rho_{\rm amp}\exp\!\left(-\frac{r^{2}}{2W_{\rm core}^{2}}\right),\qquad
\rho_{\rm amp}=\frac{f\,M_{\rm line,crit}}{2\pi W_{\rm core}^{2}},
\label{eq:ic_gauss}
\end{equation}
with $r^{2}=y^{2}+z^{2}$, $\rho_{\rm bg}=1$ in code units, $M_{\rm line,crit}=2c_s^{2}/G$ and $W_{\rm core}=0.3\,\lambda_{\rm J}$. The normalisation makes $\int_0^\infty(\rho-\rho_{\rm bg})\,2\pi r\,{\rm d}r=f\,M_{\rm line,crit}$ exactly, so the line mass converges without an imposed truncation and $f$ is set directly. The filament merges into the ambient medium, which supplies the confining pressure $P_{\rm bg}=\rho_{\rm bg}c_s^{2}$. The equilibrium-recovery control instead uses the Ostriker profile,
\begin{equation}
\rho(r)=\rho_{\rm bg}+\frac{f\rho_{c,\rm ost}}{\left[1+(r/W_{\rm core})^{2}\right]^{2}},\qquad
\rho_{c,\rm ost}=\frac{2c_s^{2}}{\pi G W_{\rm core}^{2}},
\label{eq:ic_ost}
\end{equation}
a Plummer profile of index $p=4$, whose line mass also converges. Neither uses a $p=2$ profile; the observed filaments have $p\simeq2$, which is why the width conversion of Section~\ref{sec:width_correction} is needed and why it is not invariant.

The remaining initial fields are $P(r)=\rho(r)c_s^{2}$, $v_r=v_\phi=0$, a uniform field of magnitude $B_0=c_s\sqrt{2\rho_{\rm bg}/\beta}$ at angle $\theta$ to the axis, and a longitudinal velocity seed of twelve Kolmogorov modes, $v_x=\sum_{n=1}^{12}A_n\cos(k_nx+\phi_n)$ with $k_n=2\pi n/L_x$, $A_n\propto k_n^{-11/6}$ and random phases. Domains are $L_x\times L_\perp^{2}$ with $L_x=16\,\lambda_{\rm J}$, periodic in $x$, and with the transverse boundary at $d\geq1\,\lambda_{\rm J}$ from the spine.

\smallskip\noindent\textbf{Parameters.} The filament is characterised by $f=M_{\rm line}/M_{\rm line,crit}$, the plasma $\beta=8\pi P/B^2$ and the seed-amplitude label $M_{\rm seed}=\delta v/c_s\times10^{4}$. At filament conditions ($n\sim10^4\,{\rm cm^{-3}}$, $T\sim10$\,K) $\beta=1$--$0.3$ corresponds to $B\sim15$--$30\,\mu$G \citep{Crutcher2012}. The mass-to-flux ratio scales as $\mu_\Phi/\mu_{\Phi,\rm crit}\propto f\sqrt{\beta}$ and exceeds unity throughout, so no configuration is magnetically subcritical. Table~\ref{tab:accounting} accounts for every calculation. Of the $2293$ runs in the deposited manifest, $116$ carry the numerical conclusions; the remainder are exploratory and support none of them. The exploratory set comprises the base and transition grids in $(f,\beta,\mathcal{M})$ used to locate the stable--unstable boundary, scans of the transverse boundary distance, non-ideal and perpendicular-field surveys that we show below to be numerically unresolved, and runs at $\gamma\neq1$. Those runs establish where the parameter space is well behaved; they do not measure a wavelength, and none of them is quoted as a result.

\begin{table}
\caption{Complete accounting of the campaign. Every run in the deposited manifest belongs to exactly one row. Only the first block is load-bearing: no conclusion in this paper rests on a run in the second block, and the exploratory runs are deposited so that the boundary of the well-behaved parameter space can be checked, not as measurements.}
\label{tab:accounting}
\footnotesize
\setlength{\tabcolsep}{3pt}
\begin{tabular}{@{}p{3.0cm}cp{3.5cm}@{}}
\toprule
Campaign & Runs & Conclusion supported \\
\midrule
\multicolumn{3}{@{}l}{\emph{Load-bearing (116 runs, plus the analytic helical sequence)}}\\
Equilibrium-recovery control & 4 & Section~\ref{sec:control}: pipeline recovers $22H$; contracting case gives a limit \\
Line-mass ladder (Table~\ref{tab:production}) & 10 & $\lambda_{\rm turn}^{\rm lim}$ and its resolution convergence \\
Supercritical ensemble & 39 & bead multiplicity and the ambient-confinement result \\
Contraction ladder, seed 42 & 17 & the $\rho_{\rm eff}$ scaling \\
Contraction ladder, high resolution & 4 & endpoint convergence at $512$ and $1024$ \\
Contraction-ladder replicates & 12 & realisation-to-realisation reproducibility \\
Contraction ladder, interior high resolution & 1 & convergence away from the endpoints \\
Contraction ladder, doubled transverse domain & 2 & transverse-boundary diagnostic (Section~\ref{sec:method}) \\
Contraction ladder, larger transverse domain, extended axial modes & 4 & box-size convergence of the $\rho_{\rm eff}$ relation (Section~\ref{sec:ladder}, Table~\ref{tab:bigbox}) \\
Continuous-driving campaign, $L_x=48\,\lambda_{\rm J}$ & 15 & the point-process test of Section~\ref{sec:driven}, Table~\ref{tab:driven} (9 at onset) \\
Axially modulated line mass & 8 & the inhomogeneous-assembly test of Section~\ref{sec:inhom} \\
\midrule
\multicolumn{3}{@{}l}{\emph{Exploratory (2177 runs) --- no conclusion rests on these}}\\
Base and transition grids in $(f,\beta,\mathcal{M})$ & 1560 & locate the stable--unstable boundary \\
Transverse-boundary scans & 214 & set $d\geq1\,\lambda_{\rm J}$ \\
Perpendicular-field grid and refinement ladder & 32 & shown to be Truelove-limited \\
Non-ideal (ambipolar) survey & 46 & unresolved; excluded from the argument \\
Transonic-seed, $\gamma\neq1$ and other scans & 325 & set-up robustness only \\
\midrule
\textbf{Total} & \textbf{2293} & \\
\bottomrule
\end{tabular}
\end{table}

\smallskip\noindent\textbf{Weak seeds against transonic perturbations.} The production runs are seeded once, weakly and in the linear regime, whereas real filaments are transonic. A separate set of runs replaces the linear seed with an initial transonic field at $\delta v/c_s=1$--$3$ from the same spectrum. These shorten the time to beading by a factor $3$--$5$ but change the wavelength by less than $5$ per cent, with no measurable trend against the perturbation Mach number across the ensemble, and they introduce no maximum in $\Gamma(\lambda)$ where none was present, so the plateau is not an artefact of weak seeding. The weak seed is retained for production because a growth-rate measurement requires a linear one. Continuously driven turbulence, in which energy is injected throughout the evolution rather than at $t=0$, is not tested here. The distinction is not merely technical. A continuously driven field acts as a localised, repeated compression, so it can raise the line mass in some stretches and not others while the instability is developing, which is precisely the kind of spatially varying seeding that an initial-value calculation cannot supply. It is thus a plausible route to the parsec-scale over-dispersion that our runs do not reproduce, and we identify it as an open requirement for future numerical campaigns and not as a detail of the set-up.

\smallskip\noindent\textbf{Resolution.} The load-bearing runs satisfy the Truelove criterion at the spacing-setting epoch ($N_J\approx5$--$7$ cells per local Jeans length in both directions on the isotropic grids) but sit only $1.25$--$1.75\times$ above the $N_J=4$ minimum, which is one reason we quote a bound. Truelove stability and convergence are separate requirements: $N_J>4$ protects against artificial fragmentation but does not guarantee an accurate growth rate. The trustworthy band is therefore set by the second criterion, $\Delta\Gamma/\Gamma<5$ per cent between successive resolutions, which here gives the same boundary.


\begin{table}
\caption{Load-bearing \emph{production} simulations underlying the fragmentation-wavelength result: the line-mass ladder (campaign~15) and the direct supercritical ensemble (campaign~11). Shared parameters ($f$, $\mathcal{M}$, $\beta$, $\theta$, boundary conditions and collapse/bead time) are given in each block sub-header;  These production runs, not the wider exploratory grid, carry the quoted $\lambda_{\rm turn}^{\rm lim}$.}
\label{tab:production}
\footnotesize
\setlength{\tabcolsep}{4pt}
\begin{tabular}{ccccc}
\toprule
$f$ & RNG seed & resolution & $\lambda_{\rm turn}^{\rm lim}/\lambda_{\rm J}$ & $\lambda_{\rm turn}^{\rm lim}/W$ \\
\midrule
\multicolumn{5}{@{}p{8cm}}{\emph{Line-mass ladder} (campaign~15): $\beta{=}1.0$, $\theta{=}0^\circ$, $\mathcal{M}{=}10^{-3}$ broadband seed, ambient BC (periodic/reflecting), bead $t_{\rm frag}{\approx}1.1$--$1.2\,t_{\rm J}$}\\
1.1 & 42  & $128^3$ & 1.14 & 1.95 \\
1.1 & 42  & $256^3$ & 1.14 & 1.99 \\
1.1 & 42  & $512^3$ & 1.14 & 2.00 \\
1.2 & 42  & $256^3$ & 1.00 & 1.81 \\
1.2 & 137 & $256^3$ & 0.89 & 1.60 \\
1.5 & 42  & $256^3$ & 0.89 & 1.47 \\
1.5 & 137 & $256^3$ & 0.89 & 1.47 \\
1.5 & 42  & $512^3$ & 0.89 & 1.47 \\
1.8 & 42  & $256^3$ & 0.80 & 1.19 \\
2.0 & 42  & $256^3$ & 0.80 & 1.20 \\
\midrule
\multicolumn{5}{@{}p{8cm}}{\emph{Direct supercritical ensemble} (campaign~11, 39 runs): $f{=}1.5$--$2.0$, $\beta{=}0.5$--$2.0$, $\theta{=}0^\circ$, seeds 42/137, periodic$+$reflecting ambient BC, bead $t_{\rm bead}{\approx}0.6$--$0.75\,t_{\rm J}$}\\
\multicolumn{2}{@{}c}{ensemble} & $256^3$ & $1.0$--$1.14$ & $0.89$--$1.01^{\dagger}$ \\
\bottomrule
\end{tabular}
\\[2pt]
\footnotesize $^{\dagger}$~In beam-convolved Plummer units, $\lambda/W_{\rm fil}$ (KS $p=0.17$ between periodic and reflecting BCs); the ladder's $\lambda_{\rm turn}^{\rm lim}/W$ column is in the Gaussian formation FWHM of Table~\ref{tab:production}.
\end{table}


\subsection{The reference scale}
\label{sec:reference}




For an isothermal, self-gravitating cylinder of density $\rho_0$ and sound speed $c_s$, the critical line mass for gravitational instability is \citep{Ostriker1964}:
\begin{equation}
\mu_{\rm crit} = \frac{2c_s^2}{G}
\label{eq:mucrit}
\end{equation}
The most unstable fragmentation wavelength is, as an \emph{exact} result for the isothermal equilibrium cylinder \citep{Inutsuka1992},
\begin{equation}
\lambda_{\rm IM92} = 22\,H \approx 4\times(2R_{\rm flat}),
\label{eq:lambdam}
\end{equation}
where $H = c_s\,(4\pi G\rho_c)^{-1/2}$ is the thermal scale height and $R_{\rm flat}=\sqrt8\,H$ the flattening radius (\citealt{Nagasawa1987}, eq.~33; \citealt{Fischera2012}, eqs.~5 and 24; verified independently to better than $1$ per cent in Appendix~\ref{sec:validation}), and $W_{\rm fil}\approx0.1$\,pc the characteristic filament FWHM. Expressed in FWHM units the reference depends on which FWHM is meant, and the two conventions in circulation bracket the familiar range: $\lambda_{\rm m}/W=5.07$ for the projected column-density FWHM, which is what an observer measures, and $6.04$ for the FWHM of the volume-density profile (Table~\ref{tab:benchmark}; derived in Appendix~\ref{sec:validation}). The cutoff wavelength is $\lambda_{\rm cr}=11.5\,H$, or $2.65$ in the projected convention. 

\smallskip\noindent\textbf{The reference scale under non-thermal support.} The relation above is purely thermal. Replacing $c_s$ by an effective sound speed $c_{\rm eff}^2=c_s^2+\sigma_{\rm nt}^2$ rescales every length alike, since $H\propto c_{\rm eff}$ at fixed $\rho_c$: both $\lambda_{\rm m}$ and $R_{\rm flat}$ grow by $c_{\rm eff}/c_s$, and the dimensionless ratio $\lambda_{\rm m}/W_{\rm FWHM}$ is unchanged. Transonic support, $\sigma_{\rm nt}\approx c_s$ as seen in the Taurus fibres \citep{Hacar2013}, raises both scales by $41$ per cent and leaves their ratio alone. Because we compare in units of the \emph{measured} width, the benchmark is insensitive to this rescaling at leading order; it would matter only if the non-thermal support differed systematically between the crest and the material setting the observed FWHM, which we cannot test here. Substituting $c_{\rm eff}$ treats the non-thermal motions as an isotropic pressure. That is a scaling argument and not a model of turbulence: real turbulent motions are anisotropic, correlated on the scales of interest, and can drive compression as readily as resist it.

The often-quoted ``$\lambda\approx4\times$ the width'' is $4\times$ the flat \emph{diameter} $2R_{\rm flat}$ \citep{Inutsuka1992}, which is $\approx5$--$6\times$ the \emph{FWHM}; reading it as $4\times$ the FWHM understates the scale by a third. We adopt $\lambda_{\rm m}\approx5$--$6\,W_{\rm fil}$ as the reference band, its width being the span of the two FWHM conventions and not an uncertainty. Because we compare in FWHM units and forward-model the simulations to a beam-convolved Plummer FWHM (Section~\ref{sec:width_correction}), the deficit does not depend on the equilibrium-vs-observed profile-index difference ($p=4$ vs $p\simeq2$), and we compare as the density-independent ratio $\lambda/W_{\rm fil}$. The deficit is not new in sign: \citet{Fischera2012} noted observed cores sit ``about the FWHM'' apart, and \citet{Zhang2023cal} find a spacing $\approx1.3\times$ the width in California.



Two things follow. The ratio $\lambda_{\rm m}/H=22$ and the conversion $\lambda_{\rm m}\approx4\times(2R_{\rm flat})$ are properties of the $p=4$ equilibrium and are exact within it. What is \emph{not} invariant is $\lambda_{\rm m}/W_{\rm FWHM}$, because the FWHM of a $p\simeq2$ Plummer profile is a different fraction of $H$ from that of a $p=4$ profile, and beam convolution broadens it further. The reference is thus convention-dependent at the tens-of-per-cent level before any measurement error enters, which is one reason we treat the simulated comparison through forward modelling (Section~\ref{sec:width_correction}) and not through this ratio.

With $f=\mu_{\rm line}/\mu_{\rm crit}$ and $t_{\rm frag}\propto(f-1)^{-1/2}$, marginally supercritical filaments fragment slowly. At magnetically supercritical field strength ($\beta\gtrsim0.2$) all configurations fragment for $\mathcal{M}\geq1$ at every $f$ tested once integrated adequately: above the critical line mass there is no \emph{thermal} stability boundary. This is distinct from \emph{magnetic} stabilisation, where a strong longitudinal field ($\beta\lesssim0.15$) suppresses collapse even for $f>1$ (the strong-field, magnetically supported regime; see the deposited run manifest); the thermal ($f$) and magnetic (mass-to-flux) criticalities are physically distinct.

We adopt a single reference $W_{\rm fil}=0.10$\,pc for comparability with prior HGBS work, though its universality is debated \citep{Panopoulou2017,Panopoulou2022}; the directly-measured region-Plummer widths ($0.11$--$0.16$\,pc, $\lambda/W\approx1.4$; Paper~I) show the conclusion is robust to the convention, carried as the $\sim25\%$ width systematic (Paper~I).


\begin{table}
\caption{Where the $5$--$6\times$FWHM reference comes from, for the isothermal equilibrium (Ostriker, $p=4$) cylinder. All entries are exact for that equilibrium, and are derived in Appendix~\ref{sec:validation}. The two conventions for the FWHM are the source of the quoted range: $\lambda_{\rm m}/W=5.07$ if the width is that of the projected column-density profile, as an observer measures it, and $6.04$ if it is that of the volume-density profile. Neither ratio is invariant when the observed profile has $p\simeq2$, which is why we treat the band as a reference and not a prediction.}
\label{tab:benchmark}
\footnotesize
\setlength{\tabcolsep}{4pt}
\begin{tabular}{llc}
\toprule
Quantity & Definition & Value \\
\midrule
$H$ & $c_s(4\pi G\rho_c)^{-1/2}$, scale height & $1$ \\
$R_{\rm flat}$ & flattening radius, $\sqrt{8}\,H$ & $2.828\,H$ \\
$2R_{\rm flat}$ & flat diameter & $5.657\,H$ \\
$W_{\rho}$ & FWHM of the \emph{volume}-density profile & $3.641\,H$ \\
$W_{\Sigma}$ & FWHM of the \emph{projected column}-density profile & $4.336\,H$ \\
$\lambda_{\rm cr}$ & cutoff wavelength & $11.5\,H$ \\
$\lambda_{\rm m}$ & fastest-growing wavelength & $22\,H$ \\
\midrule
$\lambda_{\rm m}/(2R_{\rm flat})$ & the ``$4\times$ diameter'' form & $3.89$ \\
$\lambda_{\rm m}/W_{\Sigma}$ & projected-FWHM convention & $5.07$ \\
$\lambda_{\rm m}/W_{\rho}$ & volume-FWHM convention & $6.04$ \\
$\lambda_{\rm cr}/W_{\Sigma}$ & cutoff, projected convention & $2.65$ \\
\bottomrule
\end{tabular}
\end{table}


\subsection{Three widths}
\label{sec:width_correction}



Three widths enter the simulation--observation comparison and must be kept separate (Fig.~\ref{fig:schematic}; Table~\ref{tab:widthdefs}): the initial Gaussian half-width $W_{\rm core}=0.3\,\lambda_{\rm J}$; the formation FWHM $W_{\rm form}=0.683\,\lambda_{\rm J}$ of the evolved projected profile ($\simeq2.28\,W_{\rm core}$, slightly below the initial $2.355\,W_{\rm core}$ owing to contraction; resolution- and seed-converged to $<1\%$); and the observed HGBS width $W_{\rm fil}$, the beam-convolved Plummer FWHM \citep{Arzoumanian2019}, which we obtain by forward-modelling the simulated filaments through the full synthetic-HGBS pipeline (projection, beam convolution, Plummer fit). This gives $\eta_W=W_{\rm form}/W_{\rm fil}=0.59$--$0.78$ and
\begin{equation}
T_1 \equiv \frac{W_{\rm form}}{W_{\rm fil}} = 0.61 \quad (\text{forward-model band } 0.59\text{--}0.78),
\label{eq:t1_correction}
\end{equation}
so that $W_{\rm fil}=W_{\rm form}/T_1\approx1.12\,\lambda_{\rm J}$. We quote the forward-model best-fit $T_1=0.61$; because it lies near the lower edge of the band, a mid-band $T_1\approx0.68$ would give a slightly \emph{larger} $\lambda/W_{\rm fil}$, so this normalisation is deliberately not the headline; the primary comparison uses the width-convention-independent growth-rate upper bound (Section~\ref{sec:control}), and this $T_1$ conversion is retained only as a cross-check. The observable is $\lambda/W_{\rm fil}$; because the fragmentation spacing $\lambda$ is a physical length independent of the width definition, the conversion from the simulation-internal ratio $\lambda/W_{\rm core}$ (normalised by the initial $\sigma=0.3\,\lambda_{\rm J}$) is
\begin{equation}
\Big(\tfrac{\lambda}{W}\Big)_{\rm fil} = \frac{W_{\rm core}}{W_{\rm fil}}\Big(\tfrac{\lambda}{W}\Big)_{\rm core} \approx 0.27\,\Big(\tfrac{\lambda}{W}\Big)_{\rm core},
\label{eq:t1_apply}
\end{equation}
i.e.\ $\lambda/W_{\rm fil}=(\lambda/\lambda_{\rm J})/1.12$; the prefactor is $0.27=W_{\rm core}/W_{\rm fil}=0.30/1.12$ (equivalently $T_1$ times the $\sigma\!\to\!$FWHM factor, $0.61\times0.44=0.27$, using $T_1\,W_{\rm core}/W_{\rm form}$), and using $T_1$ alone would overstate $\lambda/W_{\rm fil}$ by $\approx2.3\times$. This prefactor is \emph{not} a fixed constant: it carries the full $\approx30\%$ $T_1$ band ($T_1=0.59$--$0.78$), so $0.27$ is the best-fit value of a $\approx0.26$--$0.34$ range. Because the observable is sensitive to how $W_{\rm fil}$ is defined, the primary comparison uses the width-convention-independent growth-rate wavelength (an upper limit; Section~\ref{sec:control}), and this normalisation is retained only as a cross-check.



For the ambient-confined supercritical ensemble ($\lambda/W_{\rm core}=3.27$) this gives $\lambda/W_{\rm fil}\approx0.9$.


\begin{figure}
\centering
\includegraphics[width=0.99\columnwidth]{figures/fig_schematic.pdf}
\caption{Schematic of the length-scale definitions used throughout (Paper~I). \emph{(a)} Transverse column-density profile, marking the initial half-width $W_{\rm core}=\sigma_0$, the evolved spine FWHM $W_{\rm form}$, and the observed beam-convolved Plummer FWHM $W_{\rm fil}=W_{\rm form}/T_1$ ($T_1\simeq0.61$; Eq.~\ref{eq:t1_apply}). \emph{(b)} Axial view contrasting the classical widely-spaced comb ($\lambda_{\rm m}\approx5$--$6\,W$) with the observed unevenly-distributed cores: a short intra-group adjacent spacing $s_{\rm ACS,med}$ within core-rich stretches separated by core-poor filament, whose skeleton-averaged line density $L_{\rm skel}/N$ approaches the classical value under the conventions adopted in Paper~I but not under others, while the simulations give only an upper limit on any unresolved growth-rate maximum, $\lambda_{\rm turn}^{\rm lim}\lesssim1.1\,\lambda_{\rm J}$.}
\label{fig:schematic}
\end{figure}

\begin{table}
\caption{The three filament widths used in the simulation--observation comparison (all in Jeans lengths, $\lambda_{\rm J}=1$).}
\label{tab:widthdefs}
\footnotesize
\setlength{\tabcolsep}{3pt}
\begin{tabular}{p{1.5cm}p{1.4cm}p{4.0cm}}
\toprule
Symbol & Value ($\lambda_{\rm J}$) & Definition \\
\midrule
$W_{\rm core}$ & $0.30$ & initial Gaussian half-width $\sigma$ (input) \\
$W_{\rm form}$ & $0.683$ & Gaussian FWHM of evolved projected profile, supercritical ensemble ($\simeq2.28\,W_{\rm core}$) \\
$W_{\rm FWHM}$ & $0.58$ & Gaussian FWHM at the growth-rate epoch, near-critical runs (narrower, already contracting); normalises Table~\ref{tab:production} \\
$W_{\rm fil}$ & $\approx1.12$ & beam-convolved Plummer FWHM ($=W_{\rm form}/T_1$); the observable normalisation \\
\bottomrule
\end{tabular}
\end{table}

\begin{table}
\caption{Terms on the simulated wavelength $\lambda_{\rm turn}^{\rm lim}$ and on its conversion to the observable normalisation. The first block is numerical and internal to the calculations; the second enters only when a simulated wavelength is expressed in the units an observer measures, and it dominates. The corresponding terms on the observed spacing itself are given in Paper~I.}
\label{tab:simbudget}
\footnotesize
\setlength{\tabcolsep}{3pt}
\begin{tabular}{@{}p{3.1cm}p{1.9cm}p{2.3cm}@{}}
\toprule
Source & Type & Magnitude \\
\midrule
\multicolumn{3}{@{}l}{\emph{Numerical (on $\lambda_{\rm turn}^{\rm lim}$ and $\lambda_{\rm max}$)}}\\
Growth-rate resolution/epoch & modelling & $\lesssim$5\% (converged) \\
Fitting-window choice & analysis choice & $\pm1$ mode ($0.06$--$0.10$ dex) \\
Perturbation realisation & measurement & $<$2\% (4 replicates) \\
Turbulence sensitivity & measurement & $<$5\% ($N=1$ per point) \\
Line-mass ($f$) dependence & modelling & $\sim$25\% \\
\midrule
\multicolumn{3}{@{}l}{\emph{On the normalising density $\rho_{\rm eff}$}}\\
Grid refinement & numerical & $10$--$15$\% \\
Transverse box size & numerical & not converged \\
\midrule
\multicolumn{3}{@{}l}{\emph{Conversion to observable units}}\\
Gaussian vs Plummer FWHM & systematic & factor $\sim$1.5 \\
Width normalisation $\eta_W=T_1$ & systematic & $\sim$30\% \\
\bottomrule
\end{tabular}
\end{table}

The two blocks of Table~\ref{tab:simbudget} are not commensurable and we do not combine them. The numerical terms would change if the calculations were repeated at higher resolution; the conversion terms would not, because they express a fixed ambiguity in how a simulated density field is turned into the quantity an observer fits. The conversion terms are the larger, which is the reason the bound is quoted as a range, $S_{\rm sim}^{\rm upper}\lesssim1.0$--$1.3$, rather than as a single number.


\section{The controlled experiment}
\label{sec:control}


The load-bearing numerical result is a control, not a campaign. We prepared a drag-relaxed, force-balanced Ostriker cylinder and a dynamically contracting filament and passed both through the \emph{identical} measurement pipeline. For the equilibrium the pipeline recovers the classical maximum at $22H$ ($\lambda_{\rm peak}/22H=1.07$--$1.10$ at two resolutions, a genuinely resolved maximum); for the contracting filament the same pipeline finds no maximum within the trustworthy band (Fig.~\ref{fig:eqrecovery}). The difference is a property of the prepared state and not of the estimator.

Because the argument rests on that distinction, the control is shown to be stationary, not assumed to be. The cylinder is drag-relaxed to $t=6\,t_{\rm J}$ and then evolved drag-free. Over the last two Jeans times of the relaxation its gravitational binding energy changes by $0.6$ per cent and the mass-weighted radial velocity within $0.6\,\lambda_{\rm J}$ of the axis is $|v_r|_{\rm rms}\leq0.012\,c_s$; over the two Jeans times following removal of the drag, which is the interval in which growth rates are measured, the binding energy changes by $1.7$ per cent and $|v_r|_{\rm rms}\leq0.024\,c_s$. The residual drift is a slow expansion rather than a contraction, so it cannot mimic the contracting case, and if anything biases the control towards longer wavelengths, the conservative direction.

Wavelengths are measured from the growth-rate spectrum and not by counting late-time peaks. The quantity transformed is the transverse-averaged axial density contrast, $\delta(x,t)=\bar\rho(x,t)/\langle\bar\rho\rangle_x-1$ with $\bar\rho(x,t)$ the mean of $\rho$ over the two transverse directions; $P(k,t)\equiv|\hat\delta(k,t)|^2$ is its power spectrum, so $P\propto|A_k|^2$ and a fit to $\ln P$ must be halved to give an amplitude growth rate. We define $\Gamma(\lambda)\equiv\tfrac12\,{\rm d}\ln P/{\rm d}t$ throughout, so that $|A_k|\propto e^{\Gamma t}$ and $\Gamma$ is directly comparable with the linear-theory growth rate. The fit is taken over the interval in which the r.m.s.\ contrast lies between $1.5\times10^{-3}$ and $0.12$, which is linear at both ends. Seeding twelve axial modes at $\delta v/c_s=10^{-3}$ and fitting over that interval gives $\Gamma(\lambda)$ (Fig.~\ref{fig:growthrate}a), which is resolution-converged for $\lambda\gtrsim1\,\lambda_{\rm J}$ and rises with transverse resolution below $\lambda\approx0.9\,\lambda_{\rm J}$, where it is numerical.

Within the converged band $\Gamma$ plateaus and does not turn over. The statement this licenses is a logical one, and is easily overread: ${\rm d}\Gamma/{\rm d}\lambda$ does not change sign anywhere on the converged interval, so
\begin{equation}
\lambda_{\rm peak}<\lambda_{\rm min,conv}\simeq1.1\,\lambda_{\rm J}
\label{eq:bound}
\end{equation}
\emph{if a physical maximum exists at all}. We write $\lambda_{\rm turn}^{\rm lim}\equiv\lambda_{\rm min,conv}$ for that limit, and never as a measured wavelength. Where a run does show a stable, resolved maximum of $\Gamma$ we write $\lambda_{\rm max}$ instead, and we never mix the two symbols in one statement: $\lambda_{\rm turn}^{\rm lim}$ belongs to the near-critical contracting control and to the line-mass ladder, $\lambda_{\rm max}$ to the resolved runs of the impulsive-contraction ladder. "Mode selection" is likewise reserved for the second case; for the first the only defensible phrase is that the growth-rate \emph{ordering} is monotonic to the convergence limit. Forward-modelling through a synthetic beam and the same Plummer fitting used on the data (Eq.~\ref{eq:etaW}) gives the matched limit $S_{\rm sim}^{\rm upper}\lesssim1.0$--$1.3$. No preferred wavelength is detected above the numerical convergence limit. That is weaker than saying the evolving filament has no preferred wavelength: several pieces of physics omitted here could impose one below that limit, and the calculations cannot see them.

Producing a genuine maximum in an evolving cylinder requires physics this suite excludes, and it is worth naming the candidates without over-committing to any. Non-isothermal thermodynamics may eventually introduce a physical turnover in the short-wavelength growth spectrum; whether it does so at densities relevant to core formation in Gould Belt conditions requires radiative and thermal modelling beyond an isothermal calculation, and we do not attempt to specify the density at which it would act. Non-ideal effects, ambipolar diffusion in particular, damp short wavelengths preferentially once the neutral--ion coupling time exceeds the local growth time. A physical viscosity, or a turbulent cascade draining power from the coherent axial modes, would do the same by another route. Deciding among them is a separate calculation.

Two caveats temper the causal attribution. The fitted rates ($\Gamma\approx18$--$26\,t_{\rm J}^{-1}$) are the accelerating growth of a collapsing background, not quasi-static eigenvalues, so we use only the mode \emph{ordering}; and a $p=4$ initial profile gives the same order-$\lambda_{\rm J}$ scale, so the result is not a profile-index artefact.

\smallskip\noindent\textbf{What the equilibrium control does and does not validate.} It shows that the spectral estimator recovers the right answer on a stationary cylinder. It does not show that the same estimator returns an unbiased ordering on a strongly non-stationary one, and three specific concerns remain open. Radial contraction changes the axial scale during the fit, so power can move between fixed Eulerian Fourier bins without any physical mode growth; a comoving-coordinate analysis would separate the two and has not been performed. Nonlinear coupling between the twelve seeded modes could transfer power, so the broadband result is not guaranteed to reproduce what separate single-mode calculations would give. And the fitted interval, though chosen to lie in the linear regime at both ends, mixes seed growth with background amplification. Of these we have tested only the third, by shifting the interval, and it moves the answer by one axial mode (Section~\ref{sec:ladder}). Single-mode calculations at several $k$ for one equilibrium and two evolving states would settle the second and are the natural next validation step; we have not run them. The equilibrium-versus-contracting contrast is robust to the grid and box changes tested here, although single-mode and comoving-coordinate checks remain desirable validation of the detailed growth-rate ordering.






\begin{figure}
\centering
\includegraphics[width=0.88\columnwidth]{figures/fig_eqrecovery.png}
\caption{Equilibrium-recovery control. The identical growth-rate measurement pipeline is applied to a drag-relaxed, force-balanced pressure-confined Ostriker equilibrium (blue) and to a dynamically contracting filament (orange); the two differ in their prepared initial state and not in the measurement. \emph{Left:} onset fragmentation power spectrum; the equilibrium peaks at the classical $\lambda_{\rm m}\approx22H$ ($\approx5.4\,W_{\rm FWHM}$) while the contracting filament places its shortest converged scale at $\lesssim1.4\,W_{\rm FWHM}$ with no resolved maximum. \emph{Right:} $\Gamma(\lambda)$ against $\lambda/22H$; the equilibrium shows the classical peak-and-decline, the contracting filament rises monotonically. The pipeline recovers the classical wavelength when the physics is classical.}
\label{fig:eqrecovery}
\end{figure}

\begin{figure}
\includegraphics[width=0.9\columnwidth]{figures/fig_growthrate.pdf}
\caption{\emph{(a)} Growth-rate spectrum $\Gamma(\lambda)$ for a near-critical filament ($f=1.1$) at three resolutions. For $\lambda\gtrsim1\,\lambda_{\rm J}$ the curves overlap; at the grid scale (shaded) they separate and rise with resolution. $\Gamma$ plateaus rather than turning over, so $\lambda_{\rm turn}^{\rm lim}\lesssim1\,\lambda_{\rm J}$ is an \emph{upper bound}. \emph{(b)} $\lambda_{\rm turn}^{\rm lim}/W_{\rm FWHM}$ against line-mass fraction, where $W_{\rm FWHM}$ is the FWHM of the \emph{evolved, projected} simulation profile and not the initial half-width $W_{\rm core}$ or the beam-convolved observational $W_{\rm fil}$ (Table~\ref{tab:widthdefs}). The matched, beam-convolved Plummer bound is $S_{\rm sim}^{\rm upper}\lesssim1.0$--$1.3$ (grey band; the $\eta_W$ conversion is uncertain over $0.59$--$0.78$, widening it by about a quarter).}
\label{fig:growthrate}
\end{figure}


\subsection{Why only an upper limit}
\label{sec:whyupper}


Three ingredients make $\lambda_{\rm turn}^{\rm lim}$ a bound and not a measurement.

\emph{(i) Contraction competes with mode growth.} The classical calculation assumes a static, marginally-stable cylinder, so every unstable mode has unlimited time and the winner is simply the one with the largest growth rate. Our production filaments begin with $f\gtrsim1$ and no supporting radial velocity, so the spine contracts while the axial modes are still growing. The relevant comparison is between the mode growth time $\tau_{\rm grow}=\Gamma(\lambda)^{-1}$ and the time on which the background changes, $\tau_{\rm contract}=|{\rm d}\ln\rho_c/{\rm d}t|^{-1}$; long-wavelength modes have the longest growth times and are the ones least likely to complete. The argument is set out in Appendix~\ref{sec:timescale}.

The bound $S_{\rm sim}^{\rm upper}\lesssim1.0$--$1.3$ should be read as conditional on the resolution achieved: showing that a turnover has not appeared between two closely spaced resolutions is not the same as excluding one below the numerical floor.

\emph{(ii) Why no maximum is resolved.} Because the selection is set by the timescale competition rather than by the shape of $\Gamma(\lambda)$, the measured spectrum does not turn over: $\Gamma$ rises monotonically until it meets the grid scale, where the rise becomes numerical (Fig.~\ref{fig:growthrate}a). We quote the converged-band edge $\lambda_{\rm turn}^{\rm lim}$, an upper limit on where a physical turnover could lie, and never a peak.

\emph{(iii) The magnetic variables.} The field enters through $\beta$ and, equivalently, through $\mu_\Phi/\mu_{\Phi,\rm crit}\propto f\sqrt{\beta}$, which exceeds unity in every configuration run. In the uniform longitudinal-field configuration considered here a field does no work against axial modes, which compress gas along the field lines, and the simulations show no measurable dependence of the axial spectrum on $\beta$ over the tested range $0.3$--$2$ (uniform longitudinal field only; the perpendicular and helical cases are treated separately in Section~\ref{sec:optional_tests}). This is a statement about that configuration and that range, not a general result for longitudinally magnetised filaments. Its effect is radial: it resists transverse contraction and at $\beta\lesssim0.15$ halts it. A toroidal component could act on the axial mode, but only in radial force balance; seeded onto an already-contracting filament it merely pinches the spine, narrowing $W$ by $\approx25$ per cent and thereby \emph{raising} $\lambda/W$ without shortening $\lambda$.

\emph{Bridging to observable units.} Every number this paper measures is in code units, $\lambda/\lambda_{\rm J}$ or $\lambda/H$; the conversion to what an observer fits is a separate, and secondary, step. It runs in one direction only,
\begin{equation*}
\lambda\;[\lambda_{\rm J}]\;\xrightarrow{\;\div W_{\rm form}\;}\;\lambda/W_{\rm form}\;\xrightarrow{\;\times\eta_W\;}\;\lambda/W_{\rm fil}\equiv S_{\rm sim},
\end{equation*}
through the three widths of Table~\ref{tab:widthdefs}: the initial Gaussian half-width $W_{\rm core}=0.3\,\lambda_{\rm J}$, the evolved formation FWHM $W_{\rm form}=0.683\,\lambda_{\rm J}$, and the beam-convolved Plummer FWHM $W_{\rm fil}$ that an observer fits. Only the last step involves an assumption about the observation, and it is the dominant uncertainty. The bridge is one multiplicative factor,
\begin{equation}
\frac{\lambda_{\rm turn}^{\rm lim}}{W_{\rm fil}}=\eta_W\,\frac{\lambda_{\rm turn}^{\rm lim}}{W_{\rm form}},
\qquad \eta_W\equiv\frac{W_{\rm form}}{W_{\rm fil}}=T_1,
\label{eq:etaW}
\end{equation}
with $T_1=\eta_W=0.59$--$0.78$ measured by forward-modelling the simulation output through a synthetic \textit{Herschel} beam and the same Plummer fitting used on the data. The range is not a measurement error but the span of the modelling choices: it is obtained by repeating the synthetic observation at each of the four cloud distances, $135$--$436$\,pc, so that the same $18.2''$ beam corresponds to $0.012$--$0.039$\,pc of physical resolution, and by varying the fitting half-range over $0.3$--$0.5$\,pc, the background prescription between a local outer-annulus median and a global plane fit, and the Plummer index between free and fixed at $p=2$. The quoted interval is the envelope of that grid of choices and is not decomposed into independent terms, because the choices are not independent: distance and fitting half-range both act through how much of the profile wing is included. $\eta_W$ is therefore \emph{not} a universal constant of the simulation, and we do not use it as one: we quote the simulated result primarily in code units and give $S_{\rm sim}$ as a band. Every intermediate width for every load-bearing run is deposited in machine-readable form. It is Eq.~\ref{eq:etaW}, not the raw simulation number, that converts $\lambda_{\rm turn}^{\rm lim}/W_{\rm form}\approx2$ into $S_{\rm sim}^{\rm upper}\lesssim1.0$--$1.3$; the $\eta_W$ range is the dominant modelling term and widens that bound by about a quarter.




\section{The impulsive-contraction ladder}
\label{sec:ladder}


The equilibrium-versus-contracting comparison establishes the sign of the effect but not its dependence on the state of the filament, so we ran a ladder. Seventeen otherwise identical filaments were initialised from the same relaxed Ostriker cylinder with a homologous inward velocity $v_r(r)=-A\,(c_s/R_0)\,r$, for $A$ from $0$ to $1.4$. We call this a ladder in dynamical state and not in contraction rate, because the analysis below shows that $A$ is not the controlling parameter. The runs were performed at $256\times64^2$ with $\Delta x=0.0625\,\lambda_{\rm J}$ and $L_x=16\,\lambda_{\rm J}$, evolved with the same twelve-mode seed and the same growth-rate measurement, with the $A=0$ and $A=1$ endpoints repeated at $512\times128^2$.

The expectation to be tested is that mode selection follows the density at which the modes actually grow, $\lambda_{\rm select}\propto\rho_{\rm eff}^{-1/2}$, with $\rho_{\rm eff}$ the central density at the mid-point of the linear growth window. Figure~\ref{fig:ladder} shows why $A$ is not itself the controlling parameter. An imposed inward velocity does not produce steady contraction: the spine overshoots, reaching a transient central density up to $11$ times its initial value within $\approx0.3\,t_{\rm J}$, then rebounds and settles \emph{below} it, so that the most strongly kicked filaments are the least dense while their modes are growing. The relation between $A$ and $\rho_{\rm eff}$ is consequently non-monotonic, rising to $\rho_{\rm eff}\simeq3.5\,\rho_0$ near $A\approx0.2$ and falling to $0.35\,\rho_0$ at $A=1.4$. The ladder is therefore a means of sampling $\rho_{\rm eff}$ over a factor of ten and not a sequence in contraction rate.

Across that range the selected scale follows the expected scaling. Three runs, at $A=0.15$, $0.2$ and $0.25$, show no resolved maximum within the seeded modes and enter the fit as upper limits. We fit a left-censored log--log linear model: with $u_i=\ln\rho_{{\rm eff},i}$ and $y_i=\ln\lambda_{{\rm max},i}$, and writing $\mu_i=a+bu_i$ and $\sigma_i^2=\sigma_{\rm int}^2+q_i^2$, the log-likelihood is
\begin{equation}
\ln\mathcal{L}=\sum_{i\,\notin\,\mathcal{C}}\ln\varphi\!\left(\frac{y_i-\mu_i}{\sigma_i}\right)-\ln\sigma_i
+\sum_{i\,\in\,\mathcal{C}}\ln\Phi\!\left(\frac{y_i-\mu_i}{\sigma_i}\right),
\label{eq:censlik}
\end{equation}
where $\mathcal{C}$ is the censored set, $\varphi$ and $\Phi$ are the standard normal density and distribution function, $\sigma_{\rm int}$ is a fitted intrinsic scatter, and $q_i=2/(n_i\sqrt{12})$ is the half-width of the axial-quantisation bin, since only $\lambda=L_x/n$ is attainable. The likelihood is maximised by Nelder--Mead from three starting slopes and the interval on $b$ is the profile-likelihood $\Delta\ln\mathcal{L}=1/2$ contour; the implementation is deposited. Over the seventeen seed-42 runs this gives $b=-0.49\pm0.05$ with $\sigma_{\rm int}=0.09$\,dex. Discarding the censored points and fitting the remainder by least squares instead returns $-0.39$, which illustrates why the upper limits, all of which lie at the high-density end, must be included.

\smallskip\noindent\textbf{Is the relation a property of one perturbation realisation?} All seventeen runs share a single set of random phases, so the fit above is formally conditional on that realisation. We therefore repeated three rungs, chosen at the high-, intermediate- and low-density ends of the relation ($A=0.2$, $0.5$, $1.2$), with four further independent phase realisations each (Table~\ref{tab:seeds}). The result is reproducible to the resolution of the mode set. At $A=0.5$ all four replicates select $n=7$ and at $A=1.2$ all four select $n=4$, in both cases the same mode as the original, with $\lambda_{\rm max}/22H(\rho_{\rm eff})$ spanning $0.869$--$0.884$ and $0.783$--$0.789$; $\rho_{\rm eff}$ itself varies by under $4$ per cent between realisations. At $A=0.2$, where the original run is censored at $n=10$, three of four replicates are also censored at $n=10$ and the fourth resolves a maximum at $n=8$, so the censoring status at the high-density end is marginal in about a quarter of realisations. Refitting Eq.~\ref{eq:censlik} over all $29$ calculations gives
\begin{equation}
\lambda_{\rm max}\propto\rho_{\rm eff}^{-0.48\,\pm\,0.05},
\label{eq:ladderfit}
\end{equation}
with $\sigma_{\rm int}=0.08$\,dex, statistically indistinguishable from the seed-42 value and from the predicted $-1/2$ (Fig.~\ref{fig:ladder}c). The extra realisations therefore confirm the scaling without tightening it, which is the expected behaviour when the dominant uncertainty is not stochastic.

\smallskip\noindent\textbf{Does the relation survive a larger transverse domain?} It does not, and this is the principal limitation of the ladder. Section~\ref{sec:method} shows that reflecting transverse walls at $2\,\lambda_{\rm J}$ partially support the filament through image gravity, so we repeated four rungs, spanning the static, high-density, intermediate and rebounded low-density parts of the relation, in a domain twice as wide at fixed cell size, with the seeded mode set extended from twelve to twenty-four axial modes so that a shortened maximum could be resolved rather than accumulating at the top of the seeded range (Table~\ref{tab:bigbox}).

Three things change. The range of $\rho_{\rm eff}$ compresses from a factor of ten to a factor of $2.9$, because removing the image support lets every rung contract further and the strongly kicked runs no longer rebound to low density. The rank ordering breaks: the static rung now has the \emph{highest} $\rho_{\rm eff}$ and the longest selected wavelength, which is the opposite of the direction the relation requires. And a censored fit to the four runs returns a slope of $+0.15$, not $-1/2$. We therefore do not claim that $\lambda_{\rm max}\propto\rho_{\rm eff}^{-1/2}$ is established. Equation~\ref{eq:ladderfit} is a fixed-box result: within a single transverse domain the resolved modes are consistent with a gravitational scaling and reproducible between realisations, and that is all it demonstrates.

Two caveats apply to the larger-box runs themselves, so they are not a clean replacement either. Their higher densities put them closer to the Truelove floor, at $N_J\approx5$--$9$ cells per local Jeans length against $5$--$7$ in production, and two of the four select modes at $n\geq19$ of the $24$ seeded, one of them censored. Four runs also cannot constrain a slope well. The correct statement is that the fixed-box scaling is not reproduced when the boundary is moved, not that a different scaling has been measured.

What does survive, and in fact strengthens, is the distinction this paper is built on, in its ladder form: the contrast between a filament with no imposed inward motion and one that is contracting when its modes grow. In the larger box the static rung selects at $0.985$ of the instantaneous equilibrium expectation while the contracting rungs select at $0.476$--$0.649$, a cleaner separation than the $1.088$ against $0.79$--$1.12$ obtained at $L_\perp=4\,\lambda_{\rm J}$. The equilibrium-versus-contracting control is therefore robust to the transverse boundary; the quantitative ladder relation is not.

\begin{table}
\caption{Four rungs of the impulsive-contraction ladder repeated in a transverse domain twice as wide at fixed cell size ($L_\perp=8\,\lambda_{\rm J}$, $256\times128^2$), with the seeded axial mode set extended to $n\leq24$. The $A=0$ entry is the \emph{static rung of this ladder}, initialised from the relaxed Ostriker profile with no imposed inward velocity; it is not the separately prepared, drag-relaxed equilibrium-recovery control of Section~\ref{sec:control}, which was not itself repeated in the wider domain. Comparison values at the production width $L_\perp=4\,\lambda_{\rm J}$ are given in brackets. The density range compresses, the rank ordering with $\rho_{\rm eff}$ is not preserved, and a censored fit to these four runs gives $+0.15$ rather than $-1/2$. The separation between the static and the contracting rungs is preserved and widened.}
\label{tab:bigbox}
\footnotesize
\setlength{\tabcolsep}{3.5pt}
\begin{tabular}{lcccc}
\toprule
$A$ & $\rho_{\rm eff}/\rho_0$ & $n_{\rm max}$ & $\lambda_{\rm max}\,[\lambda_{\rm J}]$ & $\lambda_{\rm max}/22H(\rho_{\rm eff})$ \\
\midrule
$0.0$ (static) & $9.10$ ($1.42$) & 14 (5)  & $1.14$ ($3.20$) & $0.985$ ($1.088$) \\
$0.2$          & $7.28$ ($3.47$) & 19 (10) & $0.84$ ($<1.60$) & $0.649$ ($<0.851$) \\
$0.5$          & $4.88$ ($1.83$) & 20 (7), lim & $<0.80$ ($2.29$) & $<0.505$ ($0.882$) \\
$1.2$          & $3.14$ ($0.47$) & 17 (4)  & $0.94$ ($4.00$) & $0.476$ ($0.785$) \\
\bottomrule
\end{tabular}
\end{table}

\smallskip\noindent\textbf{How much does the fitting window matter?} More than the realisation does. Shifting the amplitude band that defines the linear window earlier or later, by factors of $0.5$ and $2.5$ on its lower edge, changes the selected mode by at most one mode number, and does so in six of the six runs tested: $n=5\to5\to6$ for $A=0$, $7\to7\to6$ for $A=0.5$, $5\to4\to4$ for $A=1.0$ and $4\to4\to3$ for $A=1.2$, with the two high-resolution endpoints behaving the same way. One mode number at $n=4$--$7$ is $0.06$--$0.10$\,dex in $\lambda$, comparable to the fitted intrinsic scatter. This is the dominant analysis-choice term, it is common to all rungs and therefore largely absorbed into the intercept rather than the slope, and it is the reason we present Eq.~\ref{eq:ladderfit} as evidence for a $\rho^{-1/2}$ scaling rather than as a precision determination of an exponent.

\begin{table}
\caption{Reproducibility of the impulsive-contraction ladder between independent perturbation realisations. Each block gives the original seed-42 run and four further realisations at the same $A$. $n_{\rm max}$ is the axial mode at which $\Gamma$ is maximal; ``lim'' marks runs in which $\Gamma$ is still rising at the shortest well-resolved seeded mode, so that the entry is an upper limit.}
\label{tab:seeds}
\footnotesize
\setlength{\tabcolsep}{4pt}
\begin{tabular}{llccc}
\toprule
$A$ & seed & $\rho_{\rm eff}/\rho_0$ & $n_{\rm max}$ & $\lambda_{\rm max}/22H(\rho_{\rm eff})$ \\
\midrule
0.20 & 42  & 3.47 & 10 (lim) & $<0.85$ \\
     & 101 & 3.39 & 10 (lim) & $<0.84$ \\
     & 202 & 3.45 & 10 (lim) & $<0.85$ \\
     & 303 & 3.50 & 10 (lim) & $<0.86$ \\
     & 404 & 3.44 & 8        & 1.06 \\
\midrule
0.50 & 42  & 1.83 & 7 & 0.882 \\
     & 101 & 1.77 & 7 & 0.869 \\
     & 202 & 1.79 & 7 & 0.874 \\
     & 303 & 1.83 & 7 & 0.884 \\
     & 404 & 1.80 & 7 & 0.875 \\
\midrule
1.20 & 42  & 0.47 & 4 & 0.785 \\
     & 101 & 0.47 & 4 & 0.783 \\
     & 202 & 0.47 & 4 & 0.785 \\
     & 303 & 0.48 & 4 & 0.789 \\
     & 404 & 0.47 & 4 & 0.784 \\
\bottomrule
\end{tabular}
\end{table}

\smallskip\noindent\textbf{Convergence of the endpoints.} Both endpoints were repeated at $512\times128^2$ and again at $1024\times256^2$, a factor of four in linear resolution and $64$ in cell count above the production grid (Table~\ref{tab:xlconv}). The selected \emph{mode} is converged: the static run peaks at $n=6$ at both $512$ and $1024$, and the dynamical run at $n=4$ at all three resolutions, so $\lambda_{\rm peak}$ takes literally the same value, $2.67$ and $4.00\,\lambda_{\rm J}$ respectively. No maximum migrates towards the grid scale as the grid is refined, which is the specific behaviour a resolution-limited spectrum would show, and no new maximum appears at shorter wavelength at $1024\times256^2$ in either case.

The residual variation in $\lambda_{\rm peak}/22H(\rho_{\rm eff})$ across resolutions, $1.09\to0.99\to1.01$ for the static run and $0.94\to0.84\to0.73$ for the dynamical one, comes not from the measured wavelength but from $\rho_{\rm eff}$, whose value at the mid-point of the fitting window depends slightly on how well the rebound is resolved. The physical statement survives at every resolution and strengthens with it: the static control returns the equilibrium result to within $1$--$9$ per cent, while the dynamical run lies $6$--$27$ per cent below the instantaneous equilibrium expectation. We quote the production-resolution median of $0.87$ in the text and note that refinement moves the dynamical endpoint further below unity, not towards it.

Convergence at the endpoints alone would not establish convergence of the fitted relation, because the endpoints are the two extremes of the density range and the slope is set by the interior. We therefore repeated one interior rung, $A=0.6$, at $512\times128^2$. It selects $n=6$, the same mode as at $256\times64^2$, so $\lambda_{\rm max}=2.67\,\lambda_{\rm J}$ is unchanged; $\rho_{\rm eff}$ falls from $1.478$ to $1.283\,\rho_0$ and $\lambda_{\rm max}/22H(\rho_{\rm eff})$ from $0.926$ to $0.863$. Four statements should be kept apart, and only the first three are converged. The discrete axial mode is converged, at the endpoints and now at an interior point. The growth-rate \emph{ordering} that selects it is converged. The measured wavelength is therefore converged. The \emph{background density trajectory} is not: $\rho_{\rm eff}$ at the mid-point of the fitting window still moves by $10$--$15$ per cent between $256$ and $512$ transverse cells, because it depends on how sharply the rebound is resolved. Since the physical inference uses the ratio $\lambda_{\rm max}/22H(\rho_{\rm eff})$, that ratio inherits the unconverged term, and this is why we quote the residual displacement as ``of order $13$ per cent'' rather than to two figures. The three interior censored runs that carry most of the leverage on the slope have not been repeated at higher resolution, and that remains an outstanding check.





\smallskip\noindent\textbf{Which hypothesis does this support?} Two must be distinguished. The first is that contraction shortens the selected wavelength directly, as a dynamical effect over and above the change in background state. The second is that the selected wavelength simply follows the instantaneous gravitational scale as the central density evolves, contraction being the means by which that density changes, not an independent wavelength-setting parameter. Within the production domain these experiments support the second much more directly than the first. An exponent of $-1/2$ is the dimensional expectation for a Jeans or scale-height problem, so recovering it would indicate that mode selection tracks $\rho_{\rm eff}$ rather than that a new mechanism is operating. The larger-domain test below shows that this reading is itself boundary-dependent, so it should be taken as an interpretation of the fixed-box ordering and not as a general result.

The point is sharpened by normalising out the changing background. Evaluating the equilibrium result at each run's own $\rho_{\rm eff}$, rather than at the initial density, removes most of the apparent shift: $\lambda_{\rm peak}/22H(\rho_{\rm eff})$ is $1.09$ for the static run and has a median of $0.87$ over the sixteen dynamical runs, spanning $0.79$--$1.12$ (Fig.~\ref{fig:ladder}d). There is a residual deficit of order $13$ per cent relative to the instantaneous equilibrium expectation, and it is consistent across the ladder, but it is an order of magnitude smaller than the factor of three to five one would infer by comparing with the wavelength of the \emph{initial} equilibrium. The correct statement is therefore that dynamical evolution permits mode selection to occur at a density different from the one defining the initial equilibrium scale, with at most a modest additional shift relative to the instantaneous scale.

Within this fixed domain the ordering follows $\rho_{\rm eff}$ rather than the imposed velocity: a filament can be given a large inward kick and still select a \emph{long} wavelength if it has already rebounded by the time its modes grow. That statement is specific to $L_\perp=4\,\lambda_{\rm J}$ and does not survive the enlargement of the transverse domain. A real filament, fed continuously rather than kicked once, need not follow the same trajectory at all.


\begin{figure*}
\includegraphics[width=0.94\textwidth]{figures/fig_ladder.pdf}
\caption{\emph{All panels refer to the production domain, $L_\perp=4\,\lambda_{\rm J}$.} The apparent $\rho_{\rm eff}^{-1/2}$ relation in panel (c) and the $0.87$ median normalisation in panel (d) are \emph{not} reproduced in the wider-domain test of Table~\ref{tab:bigbox} and must not be read as transverse-box-converged results. The impulsive-contraction ladder: seventeen runs, coloured by the imposed inward-velocity amplitude $A$. \emph{(a)} Central density against time, normalised to its initial value; the shaded band is the window over which the growth rates are fitted, and the circles mark the epoch at which $\rho_{\rm eff}$ is evaluated. The overshoot and rebound are visible directly: a strong kick raises the density transiently and leaves the filament less dense than it started while its modes are growing. \emph{(b)} Mass-weighted radial velocity within $0.6\,\lambda_{\rm J}$ of the axis, showing the same behaviour. \emph{(c)} Selected scale against $\rho_{\rm eff}$. Circles are runs with a resolved maximum of $\Gamma$; triangles are upper limits, where $\Gamma$ is still rising at the shortest well-resolved seeded mode. The solid line is the censored maximum-likelihood fit (Eq.~\ref{eq:ladderfit}) and the dashed line the $\rho_{\rm eff}^{-1/2}$ expectation. \emph{(d)} The same selected scale normalised by the equilibrium result evaluated at each run's \emph{own} $\rho_{\rm eff}$ rather than at the initial density. Once the changing background is removed the residual shift is small: the static run sits at $1.09$ and the dynamical runs have a median of $0.87$, so the deficit relative to the instantaneous equilibrium expectation is of order $13$ per cent, an order of magnitude smaller than the factor of three to five obtained by comparing with the initial equilibrium wavelength.}
\label{fig:ladder}
\end{figure*}


\section{Convergence and resolution}
\label{sec:convergence}

The claim that no maximum is resolved is a claim about a numerical spectrum, so it needs the resolution argument stated separately and in full.

\smallskip\noindent\textbf{Could numerical dissipation produce the missing maximum?} Grid-scale dissipation suppresses short-wavelength power progressively, and its signature is a short-$\lambda$ spectrum that \emph{rises} on refinement. Three results argue against a dissipative origin. The plateau edge does not move: over $128^3$--$512^3$ the fitted rate at the edge mode changes by $3.9$ and then $1.1$ per cent, while the spectrum below $\lambda\approx0.9\,\lambda_{\rm J}$ separates and rises, so the two regions behave in the opposite senses expected of a converged band and a dissipation-limited one. A cutoff at fixed $\Delta x$ would move on refinement and reveal a maximum hidden just below the band; it does not. And the axial mode set changes between resolutions, so a maximum pinned by quantisation would shift; it does not. The rows of Table~\ref{tab:production} that have been through this ladder are those at $f=1.1$ (three resolutions) and $f=1.5$ (two); the rest are single-resolution and quoted as bounds on that basis. The ladder endpoints of Section~\ref{sec:ladder} add a further, independent check at $1024\times256^2$, where the selected mode is unchanged from $512\times128^2$ and no shorter-wavelength maximum appears (Table~\ref{tab:xlconv}). These tests exclude dissipation removing a maximum lying \emph{within} the converged band; they cannot exclude one below it.

The direction in which grid-scale dissipation acts is the opposite of the one that would manufacture this result. Numerical diffusion in a second-order scheme damps preferentially at high $k$, so it \emph{suppresses} short-wavelength growth. A spectrum that is artificially smoothed near the Nyquist limit would therefore turn over, or flatten and fall, at the grid scale, and refinement would push that feature to shorter wavelengths and raise the rates behind it. What we observe below $\lambda\approx0.9\,\lambda_{\rm J}$ is the rates rising with resolution, which is the signature of a dissipation-limited region rather than a physical one, and we exclude that region from the fit for exactly that reason. Above it the rates do not change with resolution and do not turn over. The absence of a maximum in the converged band is thus not produced by numerical smoothing; if anything, smoothing would tend to create the appearance of one. What numerical dissipation does prevent is any statement about wavelengths shorter than the converged band, where a physical maximum could be hidden and where our calculations are silent.





\begin{table}
\caption{Resolution check on the two ladder endpoints. $n_{\rm peak}$ is the axial mode at which $\Gamma$ is maximal over the seeded, well-resolved modes and $\lambda_{\rm peak}=L_x/n_{\rm peak}$. The selected mode is identical at the two highest resolutions for the static run and at all three for the dynamical one; the variation in the last column follows $\rho_{\rm eff}$, not $\lambda_{\rm peak}$.}
\label{tab:xlconv}
\footnotesize
\setlength{\tabcolsep}{4pt}
\begin{tabular}{llcccc}
\toprule
Run & grid & $\rho_{\rm eff}/\rho_0$ & $n_{\rm peak}$ & $\lambda_{\rm peak}/\lambda_{\rm J}$ & $\lambda_{\rm peak}/22H$ \\
\midrule
$A=0$ & $256\times64^2$   & 0.66 & 5 & 3.20 & 1.09 \\
      & $512\times128^2$  & 0.76 & 6 & 2.67 & 0.99 \\
      & $1024\times256^2$ & 0.79 & 6 & 2.67 & 1.01 \\
\midrule
$A=1$ & $256\times64^2$   & 0.32 & 4 & 4.00 & 0.94 \\
      & $512\times128^2$  & 0.24 & 4 & 4.00 & 0.84 \\
      & $1024\times256^2$ & 0.18 & 4 & 4.00 & 0.73 \\
\bottomrule
\end{tabular}
\end{table}

\subsection{Refinement in two directions}
\label{sec:optional_tests}

Three targeted campaigns test the concerns that bear most on the interpretation, and are summarised here with the detail given in Appendix~\ref{sec:optional_app}.

Under transverse refinement ($128^3$--$512^3$) the short-wavelength part of the spectrum grows with resolution and is numerical, while the physical band $\lambda\gtrsim1\,\lambda_{\rm J}$ is converged; under axial-only refinement ($512\times128^2\to1024\times128^2$) the physical band reproduces to better than $0.5$ per cent and the short-$\lambda$ rise is unchanged, confirming it is fixed by the axial Nyquist scale and not by a physical mode migrating to shorter $\lambda$. The limit $\lambda_{\rm turn}^{\rm lim}\lesssim1.1\,\lambda_{\rm J}$ is stable in both directions.

For a perpendicular field the near-critical filament collapses monolithically in the radial direction at every resolution on a $128^3$--$512^3$ ladder, with the axial perturbation staying at seed level throughout the well-resolved phase. Short-wavelength axial structure appears only once the spine drops below $N_J\approx1$, and refinement delays its onset to deeper collapse, the signature of a grid artefact. A perpendicular field \emph{suppresses} axial fragmentation in favour of radial collapse, as static linear theory predicts.

A divergence-free helical scan ($B_\phi/B_z=0$--$2$) leaves the axial growth-rate plateau at $\lambda\approx1\,\lambda_{\rm J}$ for every winding and produces only a radial pinch that narrows $W_{\rm FWHM}$ by $\approx25$ per cent, raising $\lambda/W$ by reducing $W$ and not by shortening $\lambda$. Being seeded and not force-balanced, this configuration cannot test the confining Fiege--Pudritz sausage ($m=0$) mode \citep{Fiege2000}, which is the one geometry able to shorten $\lambda$ and remains unconstrained by any calculation presented here (Section~\ref{sec:fields}).


\section{Magnetic field geometry: what is and is not constrained}
\label{sec:fields}


The magnetic results are brief, because what they establish is mostly a limitation. A schematic oblique dispersion relation for a \emph{static} background field \citep{Nakamura1993} fixes the sign of each idealised case but cannot be used quantitatively for a magnetised cylindrical equilibrium, and in particular does not describe a current-carrying helical one; it is set out, with its restrictions, in Appendix~\ref{sec:dispersion}. This paper therefore contains no quantitative model of magnetic tension. The sign it fixes is that the magnetic term vanishes for a purely longitudinal field, reducing the problem to the field-free thermal one, and is maximal for a perpendicular field, which \emph{lengthens} the axial mode.

Table~\ref{tab:magcases} summarises the four configurations.

\begin{itemize}\setlength{\itemsep}{1pt}
\item \emph{Uniform longitudinal field}: no measurable modification of the axial scale over the tested range $\beta=0.3$--$2$.
\item \emph{Static perpendicular field}: the apparently short axial fragmentation seen at low resolution is a grid artefact, and the resolved behaviour is monolithic radial collapse, so this geometry yields no usable constraint on the axial wavelength.
\item \emph{Seeded helical field}: an out-of-equilibrium pinch that narrows $W$ without shortening $\lambda$; it is not a magnetostatic helical equilibrium and does not test one.
\item \emph{Force-balanced helical equilibrium}: not simulated. This is the only geometry identified here capable of shortening the axial mode, and it remains entirely unconstrained by this work.
\end{itemize}

Read off from the published \citet{Fiege2000} results and not computed here, a strongly toroidally dominated equilibrium shortens the $m=0$ wavelength by up to a factor $\approx2$, and our $\beta\approx0.3$--$2$ places these filaments at the weak end of that sequence. The construction a proper test would require, and the four specific obstacles to imposing such an equilibrium on a Cartesian grid, are set out in Appendix~\ref{sec:helical}.

\emph{Boundary conditions and non-ideal effects.} Periodic transverse boundaries supply an accretion channel that suppresses longitudinal fragmentation unless ambient confinement is imposed, which is why the production runs place the transverse boundary at $d\geq1\,\lambda_{\rm J}$ from the spine. At the densities probed here ($n_0\sim10^4$\,cm$^{-3}$) the ambipolar-diffusion time is long compared with the collapse time \citep{Tassis2004}, and non-ideal runs at $\eta_{\rm AD}=\{0.01,0.05\}$ still collapse radially, $12$--$15$ per cent earlier than the ideal control. Two diffusivities at one assumed density do not establish that ideal MHD is adequate across the range of ionisation fractions, field strengths and densities a real filament passes through; they establish only that within the configuration tested here the non-ideal term does not change the outcome. Because the same runs are also numerically unresolved by the Truelove criterion at fragmentation, we exclude them from the argument entirely.




\begin{table}
\caption{Schematic summary of the four magnetic-field configurations and their tested status.}
\label{tab:magcases}
\footnotesize
\setlength{\tabcolsep}{4pt}
\begin{tabular}{@{}p{2.9cm}lp{3.0cm}@{}}
\toprule
Configuration & Tested? & Effect on axial $\lambda$ \\
\midrule
Uniform longitudinal & tested & unchanged (term vanishes) \\
Static perpendicular (resolved) & tested & suppresses axial fragmentation (spindle) \\
Seeded/dynamic helical & tested & pinches $W$; does not shorten $\lambda$ \\
\midrule
\multicolumn{3}{l}{\emph{Not simulated in this work; listed for completeness}}\\
Force-balanced helical (Fiege--Pudritz) & not tested & could shorten by $\lesssim2\times$ (order-of-magnitude inference from Fiege--Pudritz, not a result of this work; Section~\ref{sec:optional_tests}) \\
\bottomrule
\end{tabular}
\end{table}


\section{Ambient confinement and the supercritical ensemble}
\label{sec:ambient}


The outcome for a supercritical filament is set by whether it is confined by an ambient medium, as real HGBS filaments are ($d$ denotes the axis-to-boundary distance in Jeans lengths). The main 654-run grid has no ambient medium, so its Gaussian skirt feeds an accretion channel through the periodic boundary that suppresses beading and drives radial collapse; with ambient confinement (\texttt{filament\_ambient} generator) the same configurations fragment longitudinally, a 12-run reconciliation suite beading at physical, $\beta$-dependent collapse times ($0.55$--$0.70\,t_{\rm J}$) that converge across boundary conditions.

Ambient-confined supercritical filaments ($f=1.5$--$2.0$, $\theta=0^\circ$) fragment longitudinally: an ensemble of 39 runs (periodic and reflecting BCs, ambient confinement) develops 3--10 interior axial density maxima at a spacing of order the Jeans length. We take the wavelength from the linear growth-rate spectrum, whose low-$n$ modes are resolution-stable: no physical turnover is detected above $\approx1\,\lambda_{\rm J}$, so any turnover is constrained only to $\lambda_{\rm turn}^{\rm lim}\lesssim1.0$--$1.14\,\lambda_{\rm J}$ (an upper bound; no resolved peak), $\lesssim0.35$--$0.5$ times the classical value (not the median peak count, which is epoch- and resolution-dependent). The ambient medium confines the filament (central-to-ambient density ratio $\sim10^2$, consistent with HGBS crest-to-background contrasts; \citealt{Arzoumanian2019}), so the adopted contrast is observationally plausible, although the calculation remains an idealised controlled experiment and not a model of an individual HGBS filament. Across the 39 runs periodic and reflecting boundaries are statistically indistinguishable (KS $p=0.17$; medians $\lambda/W_{\rm fil}\approx0.89$ and $1.01$), and this ensemble median and the growth-rate result are the same physical wavelength $\lambda\approx1.1\,\lambda_{\rm J}$ in different normalisations. The ensemble beads early ($t_{\rm bead}\approx0.6$--$0.75\,t_{\rm J}$), before runaway and at moderate contrast, so with $\lambda\approx\lambda_{\rm J}$ resolved by ${\sim}32$ cells it satisfies Truelove at measurement (Appendix~\ref{sec:validation}).

Perpendicular-field filaments ($\theta=90^\circ$) do not universally spindle in the raw low-resolution grids: an 18-run ideal-MHD grid at $d=1.0$ shows $7/18$ bead (at $\beta=0.5$ and $2.0$) with an apparently short spacing ($\lambda/W_{\rm fil}\approx0.4$), but this conflicts with static linear theory (which predicts a \emph{longer} mode) and a transverse-refinement ladder to $N_J\approx170$ resolves the near-critical case into a monolithic radial spindle, identifying the short beading as a Truelove artefact (Section~\ref{sec:optional_tests}). An exploratory ambipolar-diffusion survey was also run, but it fragments only at $N_J\approx1$--$2$, below the Truelove limit, so it cannot support any conclusion and we exclude it from the argument here; it is described, with an explicit warning that it is numerically unresolved, in the deposited material.


\section{Comparison with the observations}
\label{sec:comparison}

\subsection{What observable does a calculation of this kind predict?}


A simulation that produces one bead per wavelength predicts a \emph{line density}, one core per $\lambda$ of filament. The closest available observational proxy is the skeleton length per core, $S_{\rm skel}$ in the notation of Paper~I, and not the conditional adjacent-core spacing $S_{\rm local}$; but that proxy holds only if the measured length is fragmentation-eligible and velocity-coherent, and Paper~I shows that it is not. Neither observed statistic is therefore the theoretical observable, and no quantitative identification with either is intended here.

Two consequences are worth stating, because both cut against a naive reading. The numerical proximity of the simulated bound, $\lesssim1.0$--$1.3$ in matched units, to Paper~I's conditional value is not evidence of agreement: one is a median of observed separations conditioned on core-bearing stretches, the other a constraint on a wavelength inferred from linear growth in an idealised numerical filament, and they are not the same kind of quantity. And against the unconditional statistic the simulated bound is \emph{more} discrepant, not less, implying that fragmentation of a single coherent filament would over-produce cores per unit length unless most of the crest is core-poor, which is what Paper~I in fact measures.


\subsection{The calculations do not reproduce the observed point process}

The observed cores are not a periodic chain. Paper~I finds that they occupy only a fifth or less of the crest, that their gap distribution rejects both a periodic comb and a homogeneous Poisson process, and that no feature appears at the equilibrium-cylinder scale in the pair-correlation function. Our calculations produce $3$--$10$ near-uniformly spaced beads along a filament that is unstable along its whole length. These are different objects.

That mismatch is not a defect of the numerical set-up so much as a statement about what the set-up omits. Three ingredients are missing, and they are separable. First, the perturbation spectrum must be inhomogeneous \emph{along} the filament and not merely statistically uniform: a single realisation of a statistically homogeneous spectrum on an axially uniform filament yields approximately homogeneous fragmentation statistics, whereas the observed Fano factors of $1.6$--$25$ on parsec scales require the seed amplitude itself to vary on those scales. Continuously driven turbulence with a realistic injection scale is the natural way to supply that, and it is a different experiment from the initial-value calculations reported here. Second, the line mass must vary along the axis: a filament assembled by anisotropic accretion is supercritical in some stretches and not others, and a uniform-$f$ cylinder cannot express that. Third, the calculation must run well past first fragmentation, since migration, merging and continued accretion will erase an initially regular chain within a few free-fall times \citep{Clarke2020}.

The test of a successful calculation is accordingly not a wavelength but a set of statistics: it should reproduce the occupancy fraction as a function of scale, the Fano factor over $0.25$--$4$\,pc, the gap distribution and the pair-correlation function simultaneously. We propose those as the benchmark.

\subsection{The helical channel: the wavelength shortens but the width shrinks faster}
\label{sec:helical_result}

The one geometry identified above as able to shorten the axial mode is a force-balanced helical equilibrium. What it does to the \emph{observable} ratio is settled by the equilibrium structure itself.

We construct the equilibrium directly rather than adopting a self-similar ansatz. With the flux-frozen forms $B_z=a\rho$ and $B_\phi=b\,r\rho$, radial force balance and Poisson's equation close on
\begin{equation}
\frac{{\rm d}\rho}{{\rm d}r}\left[c_s^2+\frac{\rho\,(a^2+b^2r^2)}{4\pi}\right]
= -\rho\,\frac{{\rm d}\Phi}{{\rm d}r}-\frac{b^2r\rho^2}{2\pi},
\qquad
\frac{1}{r}\frac{{\rm d}}{{\rm d}r}\!\left(r\frac{{\rm d}\Phi}{{\rm d}r}\right)=4\pi G\rho ,
\label{eq:helix}
\end{equation}
which we integrate outward from the axis. The last term is the hoop stress: a toroidal field pinches the filament inward, so at fixed line mass it permits a denser, narrower equilibrium. That is the mechanism by which a helical field is expected to shorten the fragmentation wavelength.

It does shorten it. Holding $M_{\rm line}=M_{\rm line,crit}$ and raising the toroidal-to-poloidal energy ratio from $0$ to $150$, the central density rises by a factor of $25$ and $\lambda_{\rm m}=22H(\rho_c)$ falls by a factor of $5$, from $1.53$ to $0.31\,\lambda_{\rm J}$ (Table~\ref{tab:helix}). Taken alone that reproduces, and exceeds, the factor $\lesssim2$ read off from \citet{Fiege2000}.

But the width shrinks faster. The FWHM falls by a factor of $6.4$ over the same sequence, so the dimensionless ratio an observer measures does not fall at all: $\lambda_{\rm m}/W$ \emph{rises} monotonically from $5.93$ to $7.67$. The pinch compresses wavelength and width together and slightly favours the width. A toroidally dominated helical equilibrium therefore cannot account for the short observed $\lambda/W$ of Paper~I; it moves the prediction the wrong way.

The conclusion is robust to normalisation and to the gravity convention. Repeating the sequence at fixed \emph{central} density rather than fixed line mass gives $\lambda_{\rm m}/W$ rising from $3.11$ to $14.3$; rebuilding the equilibria against the mean-subtracted Poisson source that a periodic solver applies gives $5.95$ to $6.31$. In all three formulations the ratio rises, and in none does it fall.

\begin{table}
\caption{Force-balanced helical equilibria from Eq.~\ref{eq:helix} at fixed line mass $M_{\rm line}=M_{\rm line,crit}$, as the toroidal-to-poloidal magnetic energy ratio is raised. The absolute wavelength falls by a factor of five; the width falls faster, so the observable ratio rises. Lengths in Jeans lengths.}
\label{tab:helix}
\footnotesize
\setlength{\tabcolsep}{4pt}
\begin{tabular}{cccccc}
\toprule
$a$ & $b$ & $E_{\rm tor}/E_{\rm pol}$ & $\rho_c$ & $W_{\rm FWHM}$ & $\lambda_{\rm m}$ \; ($\lambda_{\rm m}/W$) \\
\midrule
0.35 & 0.0  & $0$    & 5.27  & 0.257 & 1.526 \; (5.93) \\
0.30 & 2.2  & $0.87$ & 6.35  & 0.232 & 1.389 \; (5.99) \\
0.22 & 3.0  & $2.1$  & 8.96  & 0.192 & 1.170 \; (6.09) \\
0.15 & 4.0  & $5.4$  & 13.80 & 0.152 & 0.943 \; (6.20) \\
0.10 & 5.5  & $14$   & 22.51 & 0.116 & 0.738 \; (6.36) \\
0.07 & 7.0  & $31$   & 35.59 & 0.089 & 0.587 \; (6.57) \\
0.05 & 9.0  & $63$   & 58.40 & 0.065 & 0.458 \; (7.01) \\
0.03 & 12.0 & $150$  & 130.3 & 0.040 & 0.307 \; (7.67) \\
\bottomrule
\end{tabular}
\end{table}

\smallskip\noindent\textbf{Status of the dynamical test.} We also implemented these equilibria as a numerical initial condition, laying the field down as the discrete curl of a vector potential on the staggered mesh so that $\nabla\!\cdot\!\mathbf{B}=0$ holds to machine precision; that removes the second of the four obstacles listed in Appendix~\ref{sec:helical}. It does not remove the third, and we state the difficulty precisely because it is the reason this channel has stayed untested. An equilibrium built against the mean-subtracted gravitational source that a periodic solver actually applies carries a true line mass well above $M_{\rm line,crit}$, and so has no stable counterpart under the unsubtracted gravity of an isolated cylinder; an equilibrium built against unsubtracted gravity is not balanced against what the solver applies. Reconciling the two needs either a genuinely isolated Poisson solve or an externally pressure-truncated transverse boundary, neither of which this set-up provides, and in test integrations the configuration consequently does not hold still. We therefore report no simulated growth rate for the helical case. The structural result above does not depend on one: Eq.~\ref{eq:helix} is an exact statement about force balance, and the behaviour of $\lambda_{\rm m}/W$ follows from the equilibrium geometry alone. A dynamical calculation could confirm the growth rate at the predicted scale; it could not reverse the direction in which the ratio moves.

\subsection{A direct test: does continuous driving produce the clustering?}
\label{sec:driven}

The natural objection to the preceding subsection is that the mismatch is an artefact of seeding the perturbations once. We therefore tested it. Athena++ supplies a continuously driven turbulence module, and we ran a campaign of fifteen supercritical filaments ($f=1.5$, Gaussian production profile, $\beta=1$) in an axial box lengthened to $L_x=48\,\lambda_{\rm J}$ at the production cell size, so that several core groups could form and be counted: three undriven controls at different perturbation seeds, and twelve driven runs spanning three injection bands (axial mode numbers $3$--$6$, $6$--$12$ and $12$--$24$, i.e.\ driving on scales above, at and below the expected fragment spacing) and two driving amplitudes, with an Ornstein--Uhlenbeck correlation time of $0.3\,t_{\rm J}$ and a solenoidal fraction of $0.6$.

Runs are compared at matched fragmentation onset rather than at matched time, because driving shortens the time to onset; we take the first snapshot at which the axial density contrast exceeds $3.5$, identify cores as connected axial stretches above twice the median density, and measure the occupancy fraction, the gap coefficient of variation and the Fano factor of core counts, the same quantities Paper~I measures on the sky. Table~\ref{tab:driven} gives the seven runs that had reached onset.

The answer is negative, and it is negative in an informative direction. Driving does increase the number of fragments, from $8$--$15$ to $16$--$18$, and roughly doubles the occupancy fraction. But it makes the resulting chain \emph{more} regular, not less: the gap coefficient of variation falls from a median of $1.50$ in the controls to $0.57$, and the Fano factor of core counts on $6\,\lambda_{\rm J}$ cells falls from $1.54$ to $0.38$, that is from mildly over-dispersed to sub-Poissonian. Neither set comes close to the observations, where Paper~I measures Fano factors of $1.6$--$25$ on parsec scales and occupancy fractions of $6$--$27$ per cent against the $1$--$3$ per cent found here.

The reason is instructive. A statistically homogeneous driver stirs every stretch of the filament equally, so it promotes fragmentation everywhere rather than in patches; it supplies additional seed power, not the persistent axial modulation of the line mass that the observed point process requires. What the observations call for is not more turbulence but an inhomogeneous \emph{assembly} history, in which some stretches become supercritical and others do not. Adding a homogeneous driver to an axially uniform filament cannot supply that, and we can now say so from calculation rather than from expectation.

\begin{table}
\caption{Continuous-driving campaign, compared at matched fragmentation onset (axial contrast $3.5$). Injection bands: L $=$ axial modes $3$--$6$, S $=$ $12$--$24$; suffix a denotes the lower driving amplitude. Cores are connected axial stretches above twice the median density in a box of $L_x=48\,\lambda_{\rm J}$. Driving raises the fragment count but \emph{lowers} the gap dispersion and the Fano factor, moving the chain away from the observed behaviour, for which Paper~I finds $F=1.6$--$25$ and $f_{\rm occ}=0.06$--$0.27$. The eight remaining driven runs had not reached onset when this analysis was made and are deposited without measurements.}
\label{tab:driven}
\footnotesize
\setlength{\tabcolsep}{4pt}
\begin{tabular}{lccccc}
\toprule
Run & $N_{\rm core}$ & $f_{\rm occ}$ & gap CoV & $F(6\,\lambda_{\rm J})$ & $F(3\,\lambda_{\rm J})$ \\
\midrule
\multicolumn{6}{@{}l}{\emph{Undriven controls}}\\
seed 42  &  8 & 0.012 & 1.34 & 1.00 & 1.50 \\
seed 137 & 14 & 0.020 & 1.50 & 1.54 & 1.41 \\
seed 271 & 15 & 0.026 & 1.60 & 1.93 & 1.53 \\
\midrule
\multicolumn{6}{@{}l}{\emph{Continuously driven}}\\
La, seed 42  & 18 & 0.031 & 0.65 & 0.08 & 0.54 \\
La, seed 137 & 17 & 0.034 & 1.44 & 0.99 & 1.47 \\
Sa, seed 42  & 16 & 0.030 & 0.50 & 0.50 & 0.38 \\
Sa, seed 137 & 16 & 0.023 & 0.40 & 0.25 & 0.25 \\
\midrule
\multicolumn{6}{@{}l}{\emph{Observed} \citep{PaperI}}\\
four clouds & --- & $0.06$--$0.27$ & $0.65$--$1.37$ & \multicolumn{2}{c}{$F=1.6$--$25$} \\
\bottomrule
\end{tabular}
\end{table}

\subsection{Testing the alternative: an axially varying line mass}
\label{sec:inhom}

The two negative results above point the same way, so we tested the alternative directly rather than leaving it as an inference. If the filament is not axially uniform, only the stretches that are locally supercritical can fragment, and the cores should appear in those stretches and nowhere else. We therefore replaced the uniform line mass with an irregular axial field,
\begin{equation}
f(x)=f_0\Big[1+\mathcal{A}\!\!\sum_{n=n_1}^{n_2}\!\cos\!\big(2\pi n x/L_x+\phi_n\big)\Big],
\label{eq:fmod}
\end{equation}
with random phases $\phi_n$ and the band $n_1$--$n_2$ chosen to give a correlation length $\ell$, normalised so that $\mathcal{A}$ sets the r.m.s.\ modulation. A single sinusoid is not adequate for this test and we do not use one: it imposes its own periodicity and simply produces one core per cycle, which we verified. Taking $f_0=1.2$, $\ell=8\,\lambda_{\rm J}$ and $\mathcal{A}=0.35$ and $0.60$ makes parts of the filament subcritical and parts strongly supercritical, and we ran four independent realisations of the modulation field at each amplitude.

The qualitative behaviour is the one required. Fragmentation is confined to the supercritical stretches: cores form in a few groups and the intervening filament, which is locally subcritical, forms none, so the empty stretches are a property of the initial line-mass field rather than of the perturbation spectrum. Neither the uniform-$f$ nor the driven calculations produce empty stretches of this kind.

The quantitative comparison is not yet a success, and the reason is specific. Measured at the same fragmentation-onset threshold used above, the modulated runs give \emph{fewer} cores than the uniform controls, $3$--$6$ against $8$--$15$, and a lower occupancy fraction, $0.004$--$0.008$ against $0.012$--$0.026$. That is a selection effect of the comparison rather than a property of the mechanism: onset is triggered by whichever patch is densest, so a strongly modulated filament reaches the threshold as soon as its single most supercritical stretch collapses, long before the others have responded. Following these runs far enough for the remaining groups to fragment requires adaptive refinement, because the first group enters a timestep runaway while the rest of the filament is still linear. We have not done that here.

\smallskip\noindent\textbf{A quantitative prediction, and a test of it.} The mechanism makes a prediction sharp enough to check against data already in hand. If fragmentation is controlled by the local line mass, the fraction of the crest that carries cores should equal the fraction that is locally supercritical: $f_{\rm occ}\simeq f_{\rm elig}$. Those two quantities are measured independently in Paper~I, the first from where the cores are and the second from the column density with the cores masked out, and they are not constructed to agree. They do agree in three clouds of four:

\begin{center}
\footnotesize
\begin{tabular}{lccc}
\toprule
Cloud & $f_{\rm occ}$ & $f_{\rm elig}$ & ratio \\
\midrule
Orion B & 0.27 & 0.19 & 1.4 \\
Perseus & 0.12 & 0.13 & 0.9 \\
Taurus  & 0.06 & 0.07 & 0.9 \\
\midrule
Aquila  & 0.27 & 0.91 & 0.3 \\
\bottomrule
\end{tabular}
\end{center}

\noindent Perseus and Taurus agree to better than $15$ per cent and Orion~B to $40$ per cent, across a factor of four in occupancy. This is the behaviour expected if the empty stretches are empty because they are subcritical, and it is not a comparison either paper set out to make. Aquila is the exception and an informative one: $91$ per cent of its crest is nominally eligible but only $27$ per cent carries a core, so there supercriticality is necessary but plainly not sufficient. Paper~I identifies Aquila as the cloud whose crest is almost entirely above the critical line mass, and as the one whose distance and line-of-sight depth are least secure.

We therefore state the conclusion as a testable proposition rather than as a preference: \emph{the occupancy fraction measures the supercritical fraction of the filament}, exactly in the calculations and to within tens of per cent in three of the four observed clouds. The way to test it properly is a dense-gas survey giving the virial line mass per stretch, which is the same measurement Paper~I already identifies as decisive.

The position is therefore this. An axially varying line mass is the only mechanism we have tested that produces core-free stretches at all, and it produces them for the right reason. Whether it also reproduces the observed Fano factors of $1.6$--$25$ is not settled here. Following every group to fragmentation within one calculation is prevented by a limitation of the method rather than of the model: a grid code following gravitational collapse without sink particles loses its timestep once the first core forms, and we verified that this is intrinsic by repeating the experiment with a stiffened equation of state, which delays the loss but does not prevent it. Sink particles, or a calculation halted and restarted group by group, would settle it.

\subsection{What the present calculations do establish}
\label{sec:whatestablished}


The result is one-sided. These calculations supply a counterexample to the assertion that a contracting filament \emph{must} fragment at the equilibrium-cylinder wavelength, and, within a fixed transverse domain, they show the selected mode following the instantaneous central density. One idealised counterexample is logically sufficient to defeat a claim of necessity; it is not sufficient to characterise the behaviour of the real population. They do not determine a preferred wavelength, and they do not reproduce the observations: the simulated filaments produce a quasi-periodic comb of 3--10 near-uniform beads, whereas the observed cores reject both periodic and Poisson descriptions and are confined to a small, clustered fraction of the crest. A successful theory must predict a point process, not merely a wavelength, and no calculation presented here does so.

One further limitation is thermodynamic. Because $\Gamma(\lambda)$ does not turn over within the converged band, an isothermal calculation contains no small-scale cutoff and cannot predict a characteristic core mass or a CMF peak. That would require $\gamma_{\rm eff}>1$ at $n\gtrsim10^{5}$--$10^{6}\,{\rm cm^{-3}}$, where dust--gas coupling and line cooling can no longer hold $T$ constant \citep{Larson1985,Glover2012}; this is a change in thermal behaviour distinct from, and at much lower density than, the classical opacity limit \citep{Rees1976,LowLyndenBell1976}. Our isothermal suite deliberately excludes it.


\section{Conclusions}
\label{sec:conclusions}

We have asked whether a dynamically evolving filament is obliged to fragment at the wavelength associated with its initial equilibrium.

\begin{enumerate}\setlength{\itemsep}{2pt}
\item \textbf{The measurement pipeline recovers the equilibrium answer when the physics is equilibrium.} Applied to a drag-relaxed, force-balanced Ostriker cylinder, whose stationarity is demonstrated rather than assumed, it returns $\lambda_{\rm peak}/22H=1.07$--$1.10$ at two resolutions and $1.01$ at a third, a genuinely resolved maximum.
\item \textbf{Applied to the near-critical contracting filament the same pipeline finds no maximum.} The growth rate rises monotonically to the edge of the converged band, so ${\rm d}\Gamma/{\rm d}\lambda$ does not change sign anywhere on that band and the only statement available is $\lambda_{\rm turn}^{\rm lim}<\lambda_{\rm min,conv}\simeq1.1\,\lambda_{\rm J}$, if a physical maximum exists at all. No preferred wavelength is detected above the numerical convergence limit; that is weaker than saying the filament has none. Taken with the previous point, this is a counterexample to the assertion that a contracting filament must fragment at the wavelength of its initial equilibrium \emph{for the idealised configuration tested here}. One counterexample defeats a claim of necessity; it does not describe what most real evolving filaments do.
\item \textbf{Within a fixed transverse domain the selected mode follows the instantaneous gravitational scale; this is not established in general.} Over $29$ calculations at $L_\perp=4\,\lambda_{\rm J}$, spanning a factor of ten in the density at which the modes grow and including four independent perturbation realisations at each of three densities, $\lambda_{\rm max}\propto\rho_{\rm eff}^{-0.48\pm0.05}$, consistent with $-1/2$, with a median displacement of $13$ per cent from the instantaneous equilibrium expectation. Both numbers are fixed-box values. Repeating four rungs in a domain twice as wide, with the seeded mode set extended to $n\leq24$, compresses the density range, breaks the rank ordering and returns a slope of $+0.15$. We therefore present the ladder as supporting evidence for the interpretation, and not as an independently established scaling law.
\item \textbf{Neither the imposed kick nor the density has been established as a general controlling variable.} Within the production box the selected mode correlates more closely with the density at the growth epoch than with the imposed kick amplitude, which overshoots and rebounds so that the most strongly kicked filaments are among the least dense while their modes grow. That correlation is not preserved when the transverse domain is enlarged, so neither $\rho_{\rm eff}$ nor the imposed velocity is demonstrated here to be a controlling variable in general.
\item \textbf{The wavelength is resolution-stable; its normalisation is not.} Both endpoints at $1024\times256^2$ and one interior rung at $512\times128^2$ select the same axial mode, and no maximum migrates in from shorter wavelength. The density at which the modes grow is less secure: it moves by $10$--$15$ per cent with grid refinement and by considerably more with transverse box size, because reflecting walls at $2\,\lambda_{\rm J}$ partially support the filament (Section~\ref{sec:method}). The ratio $\lambda/22H(\rho_{\rm eff})$ should therefore be read at the tens-of-per-cent level. Comparisons made at the same box size preserve the qualitative ordering in the doubled-box test, but the slope and normalisation of the full ladder have not been demonstrated to be transverse-box converged.
\item \textbf{The magnetic channel is now closed as an explanation.} A uniform longitudinal field leaves the axial scale unchanged over $\beta=0.3$--$2$. The ideal perpendicular-field calculation resolves monolithic radial collapse and shows the apparent short axial beading at low resolution to be numerical, so it yields no measured axial wavelength; the ambipolar-diffusion perpendicular case remains unresolved. A seeded helical field is an out-of-equilibrium pinch and not a magnetostatic equilibrium. The force-balanced helical equilibrium, previously the one geometry identified as able to shorten the axial mode, does shorten it, by a factor of five at fixed line mass; but it shrinks the filament width faster, so the observable $\lambda_{\rm m}/W$ rises rather than falls (Section~\ref{sec:helical_result}). No magnetic geometry examined here moves the predicted dimensionless spacing towards the observations.
\item \textbf{These calculations do not explain the observations.} They produce a quasi-periodic chain; the observed cores are spatially inhomogeneous and reject both a periodic comb and a homogeneous Poisson process. What a successful calculation must reproduce is a point process, and none presented here does.
\end{enumerate}

The calculations are isothermal, use idealised initial profiles, and seed the perturbations once rather than driving them continuously. Each of those omissions could introduce a physical short-wavelength cutoff that these runs cannot see. We identify one specific next calculation rather than a general appeal for more work. A statistically homogeneous initial perturbation field on an axially uniform filament tends to produce approximately homogeneous fragmentation statistics, and lacks the persistent parsec-scale modulation in instability conditions inferred in Paper~I; that is why these runs do not reproduce the observed clustering, and why tuning the parameters would not be expected to. Continuous driving breaks that symmetry, because the injected energy acts as a repeated, spatially localised compression which raises the line mass in some stretches and not others while the instability is developing. We have now carried out the driven half of that experiment (Section~\ref{sec:driven}) and it does not close the gap, which sharpens rather than removes the requirement: we have now also tested an axially \emph{varying} line mass (Section~\ref{sec:inhom}), which does produce core-free stretches; what remains is to follow it, with adaptive refinement, far enough past first fragmentation for every group to respond. The quantity to compare is not a wavelength but the occupancy fraction, the Fano factor and the gap distribution measured in Paper~I. That comparison is only meaningful once the observed network is known to be a physical object, which requires the velocity-resolved dense-gas mapping recommended in Paper~I: N$_2$H$^+$(1--0) or NH$_3$(1,1) at $\lesssim0.05$\,pc resolution. The observational and the numerical programmes therefore converge on the same missing measurement.

\section*{Acknowledgements}

We thank the \textit{Herschel} Gould Belt Survey team for the datasets, which we reprocessed. The Athena++ simulations were run on a 200 cpu cloud infrastructure (Google Cloud E2), organised from the TAURUS computer infrastructure by a Python-based pipeline run at OpenHub, Thailand (White 2026 MNRAS-RASTI submitted).

\section*{Data availability}

The supporting material, comprising the full run manifest (all $2293$ Athena++ runs with their parameters and outcomes, classified by the campaign rows of Table~\ref{tab:accounting}, of which the $116$ load-bearing runs carry the numerical result), the problem-generator source and normalisation scripts, HDF5 snapshots of the convergence subset and the equilibrium-recovery control, and the growth-rate analysis pipeline, is archived under a CC-BY-4.0 licence with the citable Zenodo DOI \texttt{10.5281/zenodo.14872318}, versioned to the accepted release. The observational material of Paper~I is deposited separately under \texttt{10.5281/zenodo.14872301}.

\nocite{PaperI}
\bibliographystyle{mnras}

\appendix

\section{A timescale interpretation}
\label{sec:timescale}

The numerical result admits a simple timescale interpretation, which we give because it is the reason we regard the counterexample as generic. It is an interpretation, not a derivation: the perturbations evolve on a time-dependent background and are not normal modes of a stationary equilibrium. In a filament of central density $\rho_c(t)$ the local Jeans length is $\lambda_{\rm J}(t)=c_s\sqrt{\pi/G\rho_c(t)}$, and a mode is realised as a fragment only if it grows before the background changes,
\begin{equation}
\tau_{\rm grow}(\lambda)\;\equiv\;\Gamma(\lambda,t)^{-1}\;\lesssim\;\tau_{\rm contract}\;\equiv\;\left|\frac{{\rm d}\ln\rho_c}{{\rm d}t}\right|^{-1}.
\label{eq:freezeout}
\end{equation}
For an equilibrium cylinder $\tau_{\rm contract}\to\infty$, every unstable mode has unlimited time, and the winner is the one with the largest $\Gamma$, the classical $\lambda_{\rm m}\approx22H$. For an evolving filament $\tau_{\rm contract}$ is of order the free-fall time while $\lambda_{\rm J}$ falls as $\rho_c^{-1/2}$, so long-wavelength modes, whose growth times are longest, are least likely to complete. The modes most likely to survive are those with $\tau_{\rm grow}\lesssim\tau_{\rm contract}$, and since $\Gamma$ increases towards shorter $\lambda$ across the converged band this selects the short end. The ladder of Section~\ref{sec:ladder} tests the resulting expectation directly. Within the production domain it exhibits the expected correlation with $\rho_{\rm eff}$, but that correlation is not reproduced when the transverse boundary is moved. The timescale argument above is therefore a plausible interpretation of the fixed-box ordering, not a demonstrated general law, which is why this appendix is labelled an interpretation and not a derivation.




\section{Resolution and field-geometry campaigns}
\label{sec:optional_app}

Three targeted campaigns test the concerns that most affect the interpretation: whether the longitudinal wavelength is resolution-limited, whether the perpendicular geometry produces genuine short beading, and whether a helical field, the one configuration expected to shorten the axial mode, actually does.

\emph{Resolution stability of $\lambda_{\rm turn}^{\rm lim}$.} Two refinement directions behave oppositely. Under \emph{transverse} refinement (the isotropic $128^3$--$512^3$ ladder) the short-wavelength spectrum ($\lambda\lesssim0.9\,\lambda_{\rm J}$, $n\gtrsim9$) grows with resolution and is numerical (unresolved transverse sub-fragmentation; Fig.~\ref{fig:growthrate}a), whereas the physical band ($\lambda\gtrsim1\,\lambda_{\rm J}$, $n\leq8$) is converged. Under \emph{axial-only} refinement ($512\times128^2\to1024\times128^2$) the physical band reproduces to better than $0.5\%$ ($\Gamma$ at $n=6,7,8$, i.e.\ $\lambda=1.33,1.14,1.00\,\lambda_{\rm J}$, is $17.6,18.7,18.7$ at both resolutions, plateauing without turning over), \emph{and} the short-$\lambda$ rise is unchanged, confirming it is fixed by the axial Nyquist scale, not a physical mode migrating to shorter $\lambda$. The $\lambda_{\rm turn}^{\rm lim}\lesssim1\,\lambda_{\rm J}$ upper bound is therefore stable: axial refinement does not move the physical-band plateau, and the only shorter-$\lambda$ structure is the transverse-resolution artefact already identified.

\emph{The perpendicular geometry: a resolved spindle, not short beading.} On a transverse-refinement ladder ($128^3$--$512^3$) the near-critical \emph{ideal} perpendicular filament collapses monolithically in the radial direction at every resolution, the axial mode staying at the box scale and the axial perturbation at seed level throughout the well-resolved phase (at $512^3$ the spine stays resolved with $N_J\approx170$ cells at $C\approx3$, falling only to $N_J\approx40$ at $C\approx30$; still $\approx10\times$ above Truelove; Appendix~\ref{sec:validation}). The axial structure at $\lambda/W\approx0.4$ emerges only after the spine has dropped below $N_J\approx1$, and each refinement postpones its emergence to a later stage of the collapse, which is the characteristic behaviour of a grid artefact and not of a physical mode. Static linear theory predicts that a perpendicular field lengthens the axial mode, and that is what the resolved calculation shows. This is for \emph{ideal} MHD; the ambipolar-perpendicular case (Section~\ref{sec:fields}) remains separately unconverged.

\emph{A seeded helical field pinches the width but does not shorten $\lambda_{\rm turn}^{\rm lim}$; the force-balanced case remains untested.} A divergence-free helical-field scan ($B_\phi/B_z=0$--$2$; construction and scan detail in the deposited material) leaves the axial growth-rate plateau at $\lambda\approx1\,\lambda_{\rm J}$ for every winding and produces only a radial \emph{pinch} that narrows $W_{\rm FWHM}$ by $\approx25\%$ (raising $\lambda/W$ by reducing $W$ and not by shortening $\lambda$; the deposited material), Being seeded and not force-balanced, this configuration cannot test the confining Fiege--Pudritz sausage ($m=0$) mode \citep{Fiege2000}, which is the one geometry able to shorten $\lambda$ and remains unconstrained by any calculation presented here (Section~\ref{sec:optional_tests}).








\section{The schematic oblique dispersion relation}
\label{sec:dispersion}

For a filament threaded by a field at an angle $\theta$ to the axis, with Alfv\'en speed $v_A$, a schematic dispersion relation valid only for a \emph{static} background field is \citep{Nakamura1993}
\begin{equation}
\omega^2 = c_s^2 k^2 + v_A^2 k^2 \sin^2\theta - 4\pi G\rho_0 F(kR),
\label{eq:dispersion}
\end{equation}
with $k$ the axial wavenumber, $\rho_0$ the central density, $R$ the filament radius and $F(kR)$ a dimensionless geometry factor from the cylindrical Poisson kernel, tending to unity as $kR\to0$ and falling monotonically to zero as $kR\to\infty$, so that self-gravity is progressively less effective at wavelengths short compared with the radius.

Equation~\ref{eq:dispersion} treats $\mathbf{B}$ as a fixed background and \emph{cannot be used quantitatively} for a magnetised cylindrical equilibrium; in particular it does not describe a current-carrying (helical) equilibrium, in which the field enters the radial force balance as well as the perturbation. We use it only as an analytic guide to the sign of each idealised case: the magnetic term $v_A^2k^2\sin^2\theta$ vanishes for a longitudinal field, so the relation reduces to the field-free thermal one, and is maximal at $\theta=90^\circ$, where it lengthens the axial mode. Consistent with the first of these, our longitudinal runs keep the growth-rate wavelength at order $\lambda_{\rm J}$ across $\beta=0.3$--$2$ with no measurable $\beta$ dependence. No quantitative model of magnetic tension is fitted or quoted anywhere in this paper.

\section{Constructing a force-balanced helical equilibrium}
\label{sec:helical}

Testing the magnetostatic helical channel needs a cylindrical MHD eigensolver, and four specific obstacles stand in the way of simply imposing such an equilibrium on a Cartesian grid. Four obstacles arise, and they are specific and not generic. The equilibrium is not available in closed form on a Cartesian grid: the Fiege--Pudritz solution is self-similar in cylindrical coordinates, so $B_z(r)$, $B_\phi(r)$ and $\rho(r)$ must be integrated numerically and then interpolated onto the mesh, and the interpolation destroys the radial force balance at the per-cent level, which is enough to launch a transient pinch on the timescale of the axial instability. Second, $\nabla\!\cdot\!\mathbf{B}=0$ must be preserved to machine precision under that interpolation, which requires the field to be laid down as a discrete curl of a vector potential consistent with the same interpolation and not as face-centred components taken from the analytic solution. Third, the configuration must be confined: a toroidally dominated equilibrium is bounded by an external pressure, so the transverse boundary condition becomes part of the equilibrium rather than an outer detail, and our ambient set-up is not constructed to supply a specified $P_{\rm ext}$. Fourth, and most restrictive, the drag-relaxation procedure we use for the Ostriker control damps \emph{all} velocities, including the ones that would reveal whether the interpolated helical state is in balance, so a different diagnostic, a residual-force map, not a velocity norm, is needed to certify stationarity before the $m=0$ mode is seeded. None of these is insurmountable, and together they define what a proper test would require.



\section{Validation of the measurement pipeline}
\label{sec:validation}


Four checks were made before the pipeline was applied; scripts and outputs are deposited.

\emph{Closed-loop injection--recovery.} Synthetic core chains of known spacing injected onto the real skeletons and recovered through the full pipeline return $\lambda/W$ to $0.4$ per cent over $\lambda/W=1.4$--$5.7$. A periodic injection over-corrects for a clustered process, which is why raw values are quoted in the main text (Paper~I).

\emph{Which FWHM the benchmark uses.} Four are in circulation: the FWHM of the volume-density profile; that of the projected column-density profile; the FWHM of a Plummer function fitted to the projected profile after background subtraction; and the beam-convolved version of the last. The classical conversion uses the second. For the Ostriker $p=4$ equilibrium, $\rho(r)=\rho_c[1+(r/R_{\rm flat})^2]^{-2}$ with $R_{\rm flat}=\sqrt8\,H$, the projected column density follows analytically,
\begin{equation}
N(x)\;\propto\;\int_{-\infty}^{\infty}\!\frac{{\rm d}y}{\left[1+(x^2+y^2)/R_{\rm flat}^2\right]^{2}}\;=\;\frac{\pi}{2}\,\frac{R_{\rm flat}}{\left[1+(x/R_{\rm flat})^{2}\right]^{3/2}},
\label{eq:ostrikercol}
\end{equation}
so the projected profile is a Plummer of index $p=3$, and its half-maximum lies where $1+(x/R_{\rm flat})^{2}=2^{2/3}$, that is at $x=R_{\rm flat}(2^{2/3}-1)^{1/2}=0.766\,R_{\rm flat}$. Hence
\begin{equation}
W_{\Sigma}=2R_{\rm flat}\left(2^{2/3}-1\right)^{1/2}=2H\sqrt{8\left(2^{2/3}-1\right)}=4.336\,H,
\label{eq:wsigma}
\end{equation}
and $\lambda_{\rm m}/W_{\Sigma}=22/4.336=5.07$. The corresponding FWHM of the \emph{volume}-density profile follows from $1+(r/R_{\rm flat})^{2}=2^{1/2}$, giving $W_{\rho}=2R_{\rm flat}(2^{1/2}-1)^{1/2}=3.641\,H$ and $\lambda_{\rm m}/W_{\rho}=6.04$. These two conventions are the origin of the quoted $5$--$6$ band, and they differ by $19$ per cent for the same physical cylinder. The observationally relevant one is $W_{\Sigma}$, since a column-density map is what is measured; we quote the band because both conventions appear in the literature. Neither ratio survives a change of profile shape: the observed filaments have $p\simeq2$, whose projection has index $1.5$, and the conversion differs.

\emph{Eigenvalue check.} The dispersion relation of the isothermal equilibrium cylinder was solved directly and agrees with \citet{Nagasawa1987} and \citet{Fischera2012} to better than $1$ per cent, peaking at $\lambda_{\rm IM92}\approx2.3\,\lambda_{\rm J}$ with a cutoff at $\lambda_{\rm cr}\approx1.2\,\lambda_{\rm J}$. Two normalisations are in use and the conversion is worth giving. Throughout this paper $\lambda_{\rm J}\equiv c_s\sqrt{\pi/G\rho_0}$ is evaluated at the initial background density $\rho_0$, whereas $H=c_s/\sqrt{4\pi G\rho_c}$ in Eq.~\ref{eq:lambdam} is evaluated at the central density $\rho_c=2.24\,\rho_0$ for the equilibrium used here. Then
\begin{equation}
\frac{22H}{\lambda_{\rm J}(\rho_0)}=\frac{22}{\sqrt{4\pi\times2.24}\times\sqrt{\pi}}=2.34,
\label{eq:hconv}
\end{equation}
which is the $2.3\,\lambda_{\rm J}$ quoted above; at a common density $\lambda_{\rm J}=2\pi H$, so the same result reads $22H=3.50\,\lambda_{\rm J}(\rho_c)$. The independent solve reproduces $22H$, and the two statements are one number in different units.

\emph{Width normalisation.} Forward-modelling the simulation output through a synthetic \textit{Herschel} beam and the same Plummer fitting used on the data gives $\eta_W=T_1=0.59$--$0.78$, with $W_{\rm form}=0.683\pm0.012\,\lambda_{\rm J}$ converged to $0.9$ per cent across resolutions (Eq.~\ref{eq:etaW}).


\bibliography{references_complete}

\bsp
\label{lastpage}
\end{document}
